\documentclass[12pt,a4paper]{report}

% =========================
% NGON NGU VA DINH DANG
% =========================
\usepackage[utf8]{inputenc}
\usepackage[T5]{fontenc}
\usepackage[vietnamese]{babel}
\usepackage[
    a4paper,
    top=2.5cm,
    bottom=2.5cm,
    left=2.5cm,
    right=2.5cm
]{geometry}
\usepackage{setspace}
\onehalfspacing
\usepackage{microtype}

% =========================
% HINH, BANG, CONG THUC
% =========================
\usepackage{graphicx}
\usepackage{float}
\usepackage{caption}
\usepackage{subcaption}
\usepackage{booktabs}
\usepackage{tabularx}
\usepackage{threeparttable}
\usepackage{longtable}
\usepackage{array}
\usepackage{adjustbox}
\usepackage{makecell}
\usepackage{ragged2e}
\usepackage{siunitx}
\usepackage{placeins}
\usepackage{enumitem}
\usepackage{amsmath}
\usepackage{amssymb}

\captionsetup{
    font=small,
    labelfont=bf,
    justification=centering
}
\sisetup{
    detect-all,
    output-decimal-marker={.},
    group-separator={,},
    group-minimum-digits=4
}
\newcolumntype{Y}{>{\RaggedRight\arraybackslash}X}
\newcolumntype{L}[1]{>{\RaggedRight\arraybackslash}p{#1}}
\emergencystretch=2em

% =========================
% SO DO VA MAU
% =========================
\usepackage{xcolor}
\usepackage{tikz}
\usetikzlibrary{arrows.meta, positioning, calc, shapes.geometric, fit, backgrounds}

\definecolor{methodblue}{RGB}{48, 92, 160}
\definecolor{methodgreen}{RGB}{52, 128, 96}
\definecolor{softgray}{RGB}{245, 247, 250}
\definecolor{linegray}{RGB}{90, 98, 110}

\tikzset{
    block/.style={
        rectangle,
        rounded corners=2pt,
        draw=linegray,
        fill=softgray,
        align=center,
        minimum height=1.0cm,
        text width=3.0cm,
        font=\small
    },
    smallblock/.style={
        rectangle,
        rounded corners=2pt,
        draw=linegray,
        fill=white,
        align=center,
        minimum height=0.78cm,
        text width=2.35cm,
        font=\footnotesize
    },
    arrow/.style={-{Latex[length=2.4mm]}, thick, draw=linegray}
}

% =========================
% LINK
% =========================
\usepackage[hidelinks]{hyperref}
\usepackage{bookmark}
\hypersetup{
    pdftitle={Phân tích cảm xúc theo khía cạnh và tóm tắt đánh giá khách sạn},
    pdfauthor={Đỗ Trần Sáng; Trần Kiết Tường},
    pdfsubject={Khóa luận tốt nghiệp},
    pdfkeywords={ABSA, opinion summarization, SemAE, HASOS, LLM judge}
}

% =========================
% THONG TIN BAO CAO
% =========================
\newcommand{\universityname}{Trường Đại học Khoa học Tự nhiên -- ĐHQG TP.HCM}
\newcommand{\facultyname}{Khoa Toán -- Tin học}
\newcommand{\advisorname}{ThS. Đoàn Thị Trâm}
\newcommand{\reporttitle}{Phân tích cảm xúc theo khía cạnh và tóm tắt đánh giá khách sạn}
\newcommand{\cityname}{TP.HCM}
\newcommand{\reportyear}{2026}

\begin{document}

% =========================
% TRANG BIA
% =========================
\begin{titlepage}
    \centering
    {\Large \universityname \par}
    \vspace{0.2cm}
    {\Large \facultyname \par}

    \vspace{2.0cm}
    {\LARGE \bfseries KHÓA LUẬN TỐT NGHIỆP ĐẠI HỌC\par}

    \vspace{1.0cm}
    {\Huge \bfseries \reporttitle \par}

    \vspace{0.8cm}
    {\Large \bfseries Tích hợp pipeline LLM/API và SemAE trên dữ liệu HASOS\par}

    \vspace{1.5cm}
    {\Large \bfseries Giảng viên hướng dẫn: \advisorname \par}

    \vspace{1.4cm}
    {\Large \bfseries Thành viên thực hiện:\par}
    \vspace{0.4cm}
    {\large
    \textbf{Đỗ Trần Sáng} \hspace{0.5cm} MSSV: 22280077\par
    \vspace{0.2cm}
    \textbf{Trần Kiết Tường} \hspace{0.5cm} MSSV: 22280102\par
    }

    \vfill
    {\large \cityname, tháng 06 năm \reportyear \par}
\end{titlepage}

\chapter*{Tóm tắt}
\addcontentsline{toc}{chapter}{Tóm tắt}

Khóa luận nghiên cứu bài toán phân tích cảm xúc theo khía cạnh (aspect-based sentiment analysis) và tóm tắt đánh giá khách sạn từ dữ liệu văn bản tự nhiên. Dữ liệu đánh giá khách sạn thường chứa nhiều nhận xét khác nhau trong cùng một văn bản: một đánh giá có thể khen vị trí nhưng chê vệ sinh hoặc dịch vụ. Do đó, tóm tắt ở cấp tổng thể không đủ để hỗ trợ người dùng hoặc nhà quản lý hiểu rõ điểm mạnh, điểm yếu theo từng khía cạnh.

Báo cáo trình bày hai hướng tiếp cận bổ trợ cho nhau. Hướng thứ nhất xây dựng pipeline dựa trên mô hình ngôn ngữ lớn thông qua API, trong đó prompt và schema JSON được dùng để trích xuất khía cạnh, cảm xúc, bằng chứng và nội dung tóm tắt. Hướng thứ hai sử dụng SemAE làm bộ chọn bằng chứng theo khía cạnh: mô hình được huấn luyện trên SPACE, sau đó suy luận trên HASOS với taxonomy 29 khía cạnh. Trên cơ sở bằng chứng được chọn, hệ thống triển khai bốn biến thể M1--M4 nhằm phân tích tác động của sinh tóm tắt trừu tượng, tách cảm xúc bằng từ khóa và tách cảm xúc bằng BERT-ABSA. Các biến thể sinh M2--M4 được thiết kế với cơ chế trung thực bằng chứng nhằm giảm thiểu hallucination, bao gồm giới hạn bằng chứng đầu vào, lời nhắc ràng buộc trung thực, bộ lọc câu không được hỗ trợ, kiểm tra nhất quán cảm xúc và cơ chế dự phòng khử trùng lặp.

Thực nghiệm trên HASOS cho thấy các biến thể sinh M3 và M4 vượt baseline trích rút M1 theo tỷ lệ vượt chuẩn (pass rate) của LLM judge: M4 đạt \num{0.787}, M3 đạt \num{0.734}, M1 đạt \num{0.720} và M2 đạt \num{0.700}. Về ROUGE, M3 đạt ROUGE-1 cao nhất (\num{0.262}), cho thấy lợi thế về độ trùng khớp từ vựng với tóm tắt tham chiếu. Kết quả này nhấn mạnh rằng hệ thống tóm tắt đánh giá không nên được đánh giá chỉ bằng ROUGE, mà cần kết hợp chỉ số tham chiếu, chỉ số vận hành và đánh giá ngữ nghĩa dựa trên bằng chứng.

\clearpage

\tableofcontents
\clearpage
\listoffigures
\clearpage
\listoftables
\clearpage

% =========================
% LOI CAM ON
% =========================
\chapter*{Lời cảm ơn}
\addcontentsline{toc}{chapter}{Lời cảm ơn}

Trong quá trình thực hiện khóa luận, chúng em đã nhận được nhiều sự hỗ trợ từ quý Thầy Cô, gia đình và bạn bè. Chúng em xin gửi lời cảm ơn chân thành đến giảng viên hướng dẫn \textbf{\advisorname} đã định hướng đề tài, góp ý phương pháp nghiên cứu, hỗ trợ quá trình thực nghiệm và giúp chúng em hoàn thiện báo cáo theo hướng có cấu trúc hơn.

Chúng em cũng xin cảm ơn quý Thầy Cô thuộc \textbf{\facultyname, \universityname} đã trang bị nền tảng kiến thức về lập trình, cơ sở dữ liệu, học máy, xử lý ngôn ngữ tự nhiên và xây dựng hệ thống phần mềm. Những kiến thức này là cơ sở quan trọng để chúng em triển khai một pipeline có khả năng xử lý dữ liệu đánh giá khách sạn, tạo đầu ra có cấu trúc và đánh giá kết quả theo nhiều góc nhìn.

Mặc dù đã cố gắng kiểm tra số liệu và tổ chức lại nội dung báo cáo, khóa luận vẫn khó tránh khỏi hạn chế. Chúng em kính mong nhận được góp ý từ quý Thầy Cô để đề tài được cải thiện trong các phiên bản tiếp theo.

\begin{flushright}
\cityname, tháng 06 năm \reportyear\\
\textbf{Sinh viên thực hiện}\\
Đỗ Trần Sáng\\
Trần Kiết Tường
\end{flushright}

\clearpage

% =========================
% LOI NOI DAU
% =========================
\chapter*{Lời nói đầu}
\addcontentsline{toc}{chapter}{Lời nói đầu}

Sự phát triển của các nền tảng du lịch trực tuyến tạo ra khối lượng lớn dữ liệu đánh giá khách sạn từ người dùng. Các đánh giá này chứa nhiều thông tin về phòng ở, vị trí, tiện nghi, vệ sinh, dịch vụ, giá cả và trải nghiệm tổng thể. Tuy nhiên, vì dữ liệu ở dạng văn bản tự nhiên, có độ dài không đồng nhất và thường chứa nhiều ý kiến trái chiều trong cùng một đánh giá, việc khai thác thủ công gặp nhiều khó khăn.

Đề tài này tập trung vào bài toán phân tích cảm xúc theo khía cạnh và tóm tắt đánh giá khách sạn. Báo cáo trình bày hai hướng tiếp cận bổ trợ cho nhau. Hướng thứ nhất xây dựng pipeline dựa trên mô hình ngôn ngữ lớn thông qua API, nhấn mạnh khả năng trích xuất thông tin có cấu trúc bằng prompt và schema JSON. Hướng thứ hai sử dụng mô hình SemAE được huấn luyện trên dữ liệu SPACE, sau đó suy luận trên tập HASOS với taxonomy 29 khía cạnh, nhằm tạo một baseline có tính tái lập cao hơn cho bước chọn bằng chứng và tóm tắt theo khía cạnh. Trên cơ sở bằng chứng do SemAE chọn, hệ thống triển khai bốn biến thể M1--M4, trong đó các biến thể sinh M2--M4 được trang bị cơ chế trung thực bằng chứng để giảm hallucination.

Báo cáo gồm bốn chương chính. Chương~1 trình bày bối cảnh, mục tiêu và đóng góp của đề tài. Chương~2 hệ thống hóa cơ sở lý thuyết về xử lý ngôn ngữ tự nhiên, phân tích cảm xúc theo khía cạnh, mô hình ngôn ngữ lớn và pipeline dữ liệu. Chương~3 mô tả hai phương pháp đề xuất, trong đó phần mở rộng tập trung vào SemAE/HASOS, các biến thể M1--M4 và cơ chế trung thực bằng chứng. Chương~4 trình bày thực nghiệm, đánh giá ROUGE, đánh giá bằng LLM judge và phân tích kết quả.

\clearpage

\chapter*{Danh sách ký hiệu}
\addcontentsline{toc}{chapter}{Danh sách ký hiệu}

Bảng~\ref{tab:symbols} liệt kê các thuật ngữ viết tắt và ký hiệu được sử dụng xuyên suốt báo cáo. Các tên riêng (ROUGE, LLM, SemAE, \dots) được giữ nguyên theo thông lệ học thuật; các thuật ngữ kỹ thuật khác được Việt hóa kèm dạng tiếng Anh trong ngoặc ở lần xuất hiện đầu tiên.

\renewcommand{\arraystretch}{1.25}
\begin{longtable}{p{0.24\textwidth} p{0.66\textwidth}}
\caption{Danh sách ký hiệu và thuật ngữ viết tắt}\label{tab:symbols}\\
\toprule
\textbf{Ký hiệu} & \textbf{Ý nghĩa} \\
\midrule
\endfirsthead
\multicolumn{2}{c}{\textit{(Tiếp theo từ trang trước)}} \\
\toprule
\textbf{Ký hiệu} & \textbf{Ý nghĩa} \\
\midrule
\endhead
\midrule
\multicolumn{2}{r}{\textit{(Tiếp tục ở trang sau)}} \\
\endfoot
\bottomrule
\endlastfoot

\multicolumn{2}{l}{\textit{\textbf{Tên riêng và viết tắt}}} \\
\addlinespace[2pt]
ABSA & Aspect-Based Sentiment Analysis, phân tích cảm xúc theo khía cạnh. \\
NLP & Natural Language Processing, xử lý ngôn ngữ tự nhiên. \\
LLM & Large Language Model, mô hình ngôn ngữ lớn. \\
BERT & Bidirectional Encoder Representations from Transformers, mô hình mã hóa ngữ cảnh hai chiều. \\
BERT-ABSA & Biến thể dùng mô hình dựa trên BERT cho nhiệm vụ phân tích cảm xúc theo khía cạnh. \\
SemAE & Semantic AutoEncoder, mô hình biểu diễn ngữ nghĩa dùng để xếp hạng câu theo khía cạnh. \\
SPACE & Tập dữ liệu đánh giá khách sạn tiếng Anh dùng để huấn luyện và đối chiếu baseline SemAE. \\
HASOS & Tập dữ liệu đánh giá khách sạn mục tiêu, có taxonomy khía cạnh chi tiết (50 khách sạn, 5.000 đánh giá, 29 khía cạnh). \\
API & Application Programming Interface, giao diện lập trình ứng dụng để gọi dịch vụ mô hình ngôn ngữ. \\
JSON & JavaScript Object Notation, định dạng dữ liệu có cấu trúc dùng để ràng buộc đầu ra của pipeline. \\
LLM judge & Mô hình ngôn ngữ lớn được dùng làm bộ chấm tự động trong thực nghiệm, theo bộ tiêu chí chấm (rubric) cố định. \\
\addlinespace[3pt]

\multicolumn{2}{l}{\textit{\textbf{Biến thể phương pháp}}} \\
\addlinespace[2pt]
M1 & Biến thể trích rút (extractive) của SemAE, giữ nguyên các câu bằng chứng được chọn. \\
M2 & Biến thể trừu tượng (abstractive) trước khi tách cảm xúc, sinh tóm tắt từ bằng chứng SemAE. \\
M3 & Biến thể tách cảm xúc bằng luật/từ khóa (keyword) trước khi sinh tóm tắt. \\
M4 & Biến thể tách cảm xúc bằng backend BERT-ABSA trước khi sinh tóm tắt. \\
\addlinespace[3pt]

\multicolumn{2}{l}{\textit{\textbf{Tham số}}} \\
\addlinespace[2pt]
$T$ & Ngưỡng chọn bằng chứng (evidence threshold); trong SemAE/HASOS, điểm là độ lệch/khoảng cách nên câu có điểm không vượt quá ngưỡng, tức $score \leq T$, được giữ làm bằng chứng. \\
$B$ & Ngân sách độ dài đầu ra; với M1 là số từ tối đa trong tóm tắt trích rút, còn với M2--M4 là số token sinh mới tối đa của mô hình sinh. \\
\addlinespace[3pt]

\multicolumn{2}{l}{\textit{\textbf{Chỉ số đánh giá và cờ lỗi}}} \\
\addlinespace[2pt]
ROUGE & Recall-Oriented Understudy for Gisting Evaluation, nhóm chỉ số so khớp n-gram giữa tóm tắt hệ thống và tóm tắt tham chiếu (báo cáo dùng ROUGE-1, ROUGE-2, ROUGE-L). \\
P@5 & Precision at 5, tỷ lệ bằng chứng đúng trong năm câu bằng chứng đầu tiên. \\
Support@5 & Tỷ lệ năm câu bằng chứng đầu thực sự hỗ trợ cho nhận định của bản tóm tắt. \\
Bằng chứng (evidence) & Câu hoặc cụm từ trong đánh giá được chọn để hỗ trợ một khía cạnh hoặc một nhận định. \\
Tỷ lệ dự phòng (fallback rate) & Tỷ lệ bản tóm tắt phải dùng cơ chế dự phòng thay vì sinh trừu tượng thành công. \\
Tỷ lệ sao chép (copy rate) & Tỷ lệ nội dung tóm tắt được sao chép trực tiếp từ câu bằng chứng. \\
Tỷ lệ nén (compression ratio) & Mức rút gọn của bản tóm tắt so với độ dài bằng chứng đầu vào. \\
Tỷ lệ vượt chuẩn (pass rate) & Tỷ lệ bản tóm tắt đạt ngưỡng chất lượng tối thiểu theo rubric của LLM judge. \\
Độ trung thực (faithfulness) & Mức độ bản tóm tắt phản ánh đúng nội dung bằng chứng, không bịa thêm thông tin. \\
Độ bao phủ (coverage) & Mức độ bản tóm tắt bao quát các ý kiến quan trọng trong bằng chứng. \\
Độ cụ thể (specificity) & Mức độ bản tóm tắt nêu chi tiết cụ thể thay vì nhận định chung chung. \\
Độ hữu ích (usefulness) & Mức độ hữu ích của bản tóm tắt đối với người đọc. \\
Khẳng định không căn cứ (unsupported claim) & Nhận định trong tóm tắt không được câu bằng chứng nào hỗ trợ. \\
Lật cảm xúc (sentiment flip) & Lỗi bản tóm tắt diễn đạt cảm xúc trái ngược với bằng chứng hoặc nhãn mục tiêu. \\
Rò rỉ khía cạnh/cảm xúc (aspect/sentiment leak) & Hiện tượng bằng chứng lẫn khía cạnh hoặc cảm xúc không thuộc mục tiêu đang xét. \\
\end{longtable}
\renewcommand{\arraystretch}{1.0}

\clearpage

% =========================
% CHUONG 1
% =========================
\chapter{Giới thiệu}

\section{Bối cảnh nghiên cứu và tính cấp thiết của đề tài}

Trong những năm gần đây, đánh giá trực tuyến trở thành một nguồn dữ liệu quan trọng phản ánh trải nghiệm thực tế của khách hàng sau khi sử dụng dịch vụ lưu trú. Một khách sạn có thể nhận hàng trăm đến hàng nghìn đánh giá, trong đó mỗi đánh giá thường đề cập đồng thời nhiều khía cạnh như vị trí, phòng, nhân viên, bữa sáng, tiện nghi, vệ sinh hoặc giá cả. Nếu chỉ đọc điểm sao trung bình, nhà quản lý khó nhận biết nguyên nhân cụ thể tạo nên sự hài lòng hoặc không hài lòng của khách hàng. Bài toán khai thác và tóm tắt đánh giá khách hàng đã được nghiên cứu từ sớm trong hướng review mining~\cite{hu2004mining}.

Việc phân tích tự động dữ liệu đánh giá khách sạn có ý nghĩa thực tiễn rõ rệt. Hệ thống có thể giúp phát hiện các vấn đề lặp lại, theo dõi thay đổi chất lượng dịch vụ theo thời gian, so sánh các chi nhánh và tạo báo cáo ngắn gọn cho người quản lý. Ở chiều ngược lại, người dùng mới cũng có thể được hỗ trợ bởi các bản tóm tắt tập trung vào những khía cạnh họ quan tâm thay vì phải đọc nhiều đánh giá dài.

Tuy nhiên, dữ liệu đánh giá có nhiều đặc điểm gây khó cho các phương pháp xử lý tự động. Câu đánh giá thường ngắn hoặc thiếu cấu trúc, có cách diễn đạt đa dạng, chứa tiếng lóng, lỗi chính tả hoặc nhiều ý kiến trong cùng một câu. Một đánh giá có thể tích cực về vị trí nhưng tiêu cực về tiếng ồn, hoặc khen nhân viên nhưng phàn nàn về vệ sinh. Do đó, gán một nhãn cảm xúc chung cho toàn bộ đánh giá sẽ làm mất thông tin chi tiết ở cấp khía cạnh, vốn là trọng tâm của các benchmark ABSA như SemEval-2014 Task 4~\cite{pontiki2014semeval}.

Sự xuất hiện của các mô hình ngôn ngữ lớn mở ra một hướng tiếp cận linh hoạt cho bài toán này. Thông qua prompt và schema phù hợp, LLM có thể trích xuất khía cạnh, xác định cảm xúc, chọn bằng chứng và tạo tóm tắt mà không cần huấn luyện từ đầu. Dù vậy, pipeline LLM/API cần kiểm soát chi phí, độ trễ, lỗi định dạng, khả năng chạy lại và tính nhất quán của kết quả. Đây là lý do báo cáo không chỉ xem xét một pipeline LLM trực tiếp, mà còn bổ sung phương pháp thứ hai dựa trên SemAE để tách rõ bước chọn bằng chứng khỏi bước sinh văn bản.

\section{Mục tiêu và phạm vi nghiên cứu}

Mục tiêu tổng quát của đề tài là xây dựng và đánh giá pipeline phân tích cảm xúc theo khía cạnh và tóm tắt đánh giá khách sạn. Cụ thể, đề tài hướng tới các mục tiêu sau:

\begin{itemize}[leftmargin=1.2cm]
    \item Chuẩn hóa dữ liệu đánh giá khách sạn thành các đơn vị xử lý phù hợp cho bài toán ABSA và tóm tắt.
    \item Thiết kế pipeline LLM/API có schema đầu ra rõ ràng, có kiểm tra định dạng, cache và cơ chế xử lý lỗi.
    \item Triển khai phương pháp SemAE trên dữ liệu HASOS để chọn bằng chứng theo 29 khía cạnh khách sạn.
    \item So sánh các biến thể biểu diễn đầu ra gồm trích rút (extractive), trừu tượng (abstractive), tách cảm xúc bằng từ khóa (keyword) và tách cảm xúc bằng BERT-ABSA.
    \item Thiết kế cơ chế trung thực bằng chứng cho các biến thể sinh nhằm giảm hallucination và cải thiện độ tin cậy theo LLM judge.
    \item Đánh giá kết quả bằng ROUGE, chỉ số vận hành và LLM judge theo bộ tiêu chí chấm (rubric) cố định.
\end{itemize}

Phạm vi nghiên cứu tập trung vào dữ liệu đánh giá khách sạn dạng văn bản. Báo cáo không huấn luyện LLM từ đầu, không xây dựng hạ tầng phân tán quy mô sản xuất và không báo các chỉ số yêu cầu nhãn gold không tồn tại trong dữ liệu hiện có, chẳng hạn ABSA Macro-F1 ở cấp câu hoặc Quad F1 ở cấp span.

\section{Phương pháp tiếp cận và đóng góp của đề tài}

Đề tài được tổ chức quanh hai phương pháp. Phương pháp thứ nhất dùng LLM thông qua API để phân tích đánh giá theo prompt và schema có cấu trúc. Phương pháp này phù hợp khi cần linh hoạt về định dạng đầu ra và muốn tận dụng năng lực suy luận ngôn ngữ của mô hình lớn. Phương pháp thứ hai dùng SemAE để xếp hạng câu theo khía cạnh, sau đó tạo các bản tóm tắt ở nhiều dạng khác nhau. Phương pháp này giúp quá trình chọn bằng chứng có khả năng tái lập tốt hơn, đồng thời tạo điều kiện phân tích ảnh hưởng của việc tách cảm xúc (sentiment split) và token budget.

Đóng góp chính của đề tài gồm:

\begin{itemize}[leftmargin=1.2cm]
    \item Một thiết kế pipeline hai nhánh cho bài toán phân tích và tóm tắt theo khía cạnh trong miền khách sạn.
    \item Một quy trình chuyển SemAE từ dữ liệu huấn luyện SPACE sang suy luận trên HASOS với taxonomy 29 khía cạnh.
    \item Một bộ biến thể M1--M4 để tách riêng ảnh hưởng của chọn bằng chứng, sinh tóm tắt và tách cảm xúc.
    \item Một cơ chế trung thực bằng chứng gồm năm thành phần (giới hạn bằng chứng đầu vào, lời nhắc ràng buộc, bộ lọc câu không được hỗ trợ, kiểm tra nhất quán cảm xúc và dự phòng khử trùng lặp) giúp các biến thể sinh giảm hallucination và vượt baseline trích rút theo LLM judge.
    \item Một bộ đánh giá kết hợp ROUGE, sweep tham số và LLM judge, có ghi rõ chỉ số nào đủ điều kiện báo cáo chính thức.
\end{itemize}

% =========================
% CHUONG 2
% =========================
\chapter{Cơ sở lý thuyết}

\section{Xử lý ngôn ngữ tự nhiên}

Xử lý ngôn ngữ tự nhiên (NLP) là lĩnh vực nghiên cứu các phương pháp giúp máy tính phân tích, biểu diễn và sinh ngôn ngữ con người. Trong bài toán đánh giá khách sạn, NLP được dùng để tách câu, làm sạch văn bản, biểu diễn câu thành vector ngữ nghĩa, phát hiện khía cạnh được nhắc đến và sinh bản tóm tắt ngắn gọn.

Một pipeline NLP thường bắt đầu bằng tiền xử lý văn bản. Các bước phổ biến gồm chuẩn hóa khoảng trắng, tách câu, loại bỏ ký tự lỗi, kiểm tra ngôn ngữ và chuẩn hóa đơn vị đầu vào. Với dữ liệu đánh giá, bước tách câu đặc biệt quan trọng vì một đánh giá dài có thể chứa nhiều nhận xét khác nhau, trong khi phương pháp SemAE sử dụng câu làm đơn vị xếp hạng bằng chứng.

\section{Phân tích cảm xúc và ABSA}

Phân tích cảm xúc (sentiment analysis) nhằm xác định thái độ tích cực, tiêu cực hoặc trung tính trong văn bản. Ở cấp tài liệu (document-level), hệ thống đưa ra một nhãn cho toàn bộ đánh giá. Ở cấp câu (sentence-level), nhãn được gán cho từng câu. Tuy nhiên, trong miền khách sạn, hai cấp này vẫn chưa đủ vì một câu có thể chứa nhiều khía cạnh.

Phân tích cảm xúc theo khía cạnh (ABSA) giải quyết hạn chế trên bằng cách gắn cảm xúc với từng khía cạnh cụ thể~\cite{pontiki2014semeval}. Ví dụ, câu ``phòng sạch nhưng bữa sáng nghèo nàn'' chứa cảm xúc tích cực cho khía cạnh phòng và cảm xúc tiêu cực cho khía cạnh bữa sáng. Trong báo cáo này, ABSA đóng vai trò cầu nối giữa đánh giá thô và tóm tắt theo từng nhóm thông tin.

\section{Mô hình ngôn ngữ lớn và prompt có cấu trúc}

Mô hình ngôn ngữ lớn (LLM) có khả năng hiểu ngữ cảnh và sinh văn bản theo yêu cầu. Phần lớn các mô hình hiện nay dựa trên kiến trúc Transformer với cơ chế tự chú ý (self-attention)~\cite{vaswani2017attention}; các mô hình tiền huấn luyện như BERT cũng là nền tảng quan trọng cho nhiều backend phân tích cảm xúc~\cite{devlin2019bert}. Khi được kết hợp với kỹ thuật thiết kế prompt (prompt engineering), LLM có thể thực hiện các tác vụ như trích xuất khía cạnh, chọn bằng chứng, đánh giá độ trung thực và tạo tóm tắt~\cite{brown2020language,bhaskar2023prompted}. Để giảm rủi ro đầu ra không ổn định, pipeline cần yêu cầu schema rõ ràng, kiểm tra JSON, thử lại (retry) khi lỗi và lưu bộ nhớ đệm (cache) kết quả.

Trong phương pháp thứ nhất, LLM được dùng trực tiếp như thành phần phân tích chính. Trong phần đánh giá của phương pháp thứ hai, LLM được dùng như bộ chấm (judge) để chấm các yếu tố khó đo bằng ROUGE, chẳng hạn mức hỗ trợ của bằng chứng, lật cảm xúc (sentiment flip) hoặc sai khía cạnh~\cite{liu2023geval,zheng2023judging}. Hai vai trò này cần được phân biệt: LLM sinh kết quả là một phần của hệ thống, còn LLM judge là công cụ đánh giá theo bộ tiêu chí chấm.

\section{SemAE và tóm tắt theo khía cạnh}

SemAE là hướng tiếp cận dùng biểu diễn ngữ nghĩa để chọn câu liên quan đến một khía cạnh~\cite{basuroychowdhury2022sparse}. Thay vì yêu cầu LLM đọc toàn bộ đánh giá, SemAE mã hóa câu đánh giá và so sánh với prototype của từng khía cạnh. Các câu có điểm gần với prototype được xếp hạng cao hơn và được dùng làm bằng chứng cho tóm tắt.

SemAE được huấn luyện trên dữ liệu SPACE, một tập đánh giá khách sạn tiếng Anh có tóm tắt tham chiếu (gold summary) theo các khía cạnh tổng quát~\cite{angelidis2021extractive}. Khi chuyển sang HASOS, hệ thống dùng taxonomy 29 khía cạnh và bộ từ khóa hạt giống (seed words) riêng cho miền khách sạn. Thiết kế này tách bước học biểu diễn ngữ nghĩa khỏi bước định nghĩa taxonomy mục tiêu.

\section{Đánh giá tóm tắt}

ROUGE là nhóm chỉ số phổ biến để so sánh tóm tắt hệ thống với tóm tắt tham chiếu (gold summary) dựa trên n-gram~\cite{lin2004rouge}. ROUGE hữu ích để đo mức độ trùng khớp từ vựng, nhưng không đủ để đánh giá đầy đủ độ trung thực (faithfulness), đúng khía cạnh hoặc đúng cảm xúc. Vì vậy, báo cáo kết hợp ROUGE với bộ chấm tự động bằng mô hình ngôn ngữ lớn (LLM judge) theo bộ tiêu chí chấm (rubric) cố định~\cite{liu2023geval,zheng2023judging}. Cách đánh giá này giúp phân biệt hai câu hỏi: bản tóm tắt có giống tóm tắt tham chiếu không, và bản tóm tắt có được bằng chứng hỗ trợ không.

% =========================
% CHUONG 3
% =========================
\chapter{Phương pháp đề xuất}

\section{Tổng quan thiết kế hai phương pháp}

Báo cáo giữ phương pháp LLM/API đã trình bày trong bản gốc và bổ sung phương pháp SemAE/HASOS như một nhánh thực nghiệm thứ hai. Hai phương pháp cùng hướng tới đầu ra có cấu trúc theo khía cạnh, nhưng khác nhau ở cách chọn bằng chứng và mức độ phụ thuộc vào LLM.

\begin{figure}[tbp]
\centering
\begin{tikzpicture}[node distance=0.9cm and 0.8cm]
    \node[block, text width=3.2cm] (reviews) {Đánh giá khách sạn};
    \node[block, below left=of reviews, draw=methodblue] (m1) {Phương pháp 1\\LLM/API};
    \node[block, below right=of reviews, draw=methodgreen] (m2) {Phương pháp 2\\SemAE/HASOS};
    \node[block, below=of m1, text width=3.6cm] (struct)
{Tóm tắt theo\\cluster evidence\\trong từng aspect/sentiment}
    \node[block, below=of m2] (summ) {Tóm tắt theo\\khía cạnh};
    \coordinate (mid) at ($(struct)!0.5!(summ)$);
    \node[block, below=of mid, text width=3.4cm] (eval) {Đánh giá\\ROUGE + LLM judge};
    \draw[arrow] (reviews) -| (m1);
    \draw[arrow] (reviews) -| (m2);
    \draw[arrow] (m1) -- (struct);
    \draw[arrow] (m2) -- (summ);
    \draw[arrow] (struct) |- (eval);
    \draw[arrow] (summ) |- (eval);
\end{tikzpicture}
\caption{Thiết kế tổng quan với hai phương pháp bổ trợ cho nhau}
\label{fig:two-method-overview}
\end{figure}

Hình~\ref{fig:two-method-overview} cho thấy hai phương pháp được thiết kế như hai nhánh bổ trợ. Phương pháp~1 phù hợp với bài toán cần khả năng diễn giải linh hoạt và đầu ra JSON trực tiếp. Phương pháp~2 phù hợp với mục tiêu nghiên cứu có tính tái lập: cùng một bộ bằng chứng được chọn bởi SemAE có thể được biểu diễn theo nhiều cách khác nhau, từ đó phân tích tác động của tách cảm xúc và sinh tóm tắt.

\section{Phương pháp 1: Quy trình LLM/API}

Phương pháp~1 sử dụng mô hình ngôn ngữ lớn thông qua API để xây dựng một quy trình phân tích và tóm tắt đánh giá khách sạn dựa trên khía cạnh. Trong phạm vi báo cáo này, phương pháp được trình bày theo hướng xử lý ở mức câu hoặc đơn vị ý kiến, tức là hệ thống không gán một nhãn chung cho toàn bộ đánh giá, mà tách đánh giá thành các đoạn nhỏ có thể truy vết, sau đó xác định khía cạnh, cảm xúc, cụm chủ đề và bằng chứng cho từng đoạn. Cách tiếp cận này phù hợp với bản chất của đánh giá khách sạn, nơi một đoạn phản hồi ngắn cũng có thể chứa nhiều ý kiến khác nhau về phòng, tiện ích, nhân viên, vị trí, trải nghiệm tổng thể hoặc ý định quay lại.

Khác với cách yêu cầu LLM sinh trực tiếp một đoạn tóm tắt từ toàn bộ dữ liệu thô, phương pháp này xem tóm tắt là tầng cuối của một chuỗi xử lý có cấu trúc. Đánh giá đầu vào trước hết được đọc theo lô, làm sạch và gắn định danh khách sạn. Sau đó, văn bản được tách thành các đơn vị nhỏ hơn, mỗi đơn vị được đưa qua bước trích xuất ABSA để xác định khía cạnh và đoạn bằng chứng. Kết quả khía cạnh không được dùng một cách tuyệt đối mà tiếp tục được kiểm soát bằng luật ranh giới giữa các nhãn. Tiếp theo, hệ thống phân loại cảm xúc theo từng đoạn, gán cụm chủ đề trong từng khía cạnh và gom các bằng chứng cùng cụm để sinh tóm tắt. Do đó, tóm tắt cuối không phải là sản phẩm của một lần sinh văn bản tự do, mà là kết quả của quá trình tích lũy và tổ chức bằng chứng.

Điểm chính của phương pháp nằm ở sự kết hợp giữa năng lực hiểu ngữ nghĩa của LLM và các cơ chế kiểm soát của hệ thống dữ liệu. LLM được dùng ở những bước cần hiểu văn bản tự nhiên, chẳng hạn tách ý kiến trong câu dài, nhận diện khía cạnh, chuẩn hóa đoạn văn song ngữ, phân loại cảm xúc và sinh tóm tắt. Ngược lại, các thành phần như đọc dữ liệu theo lô, bộ nhớ đệm, điểm lưu tiến độ, kiểm tra phản hồi JSON, thử lại khi lỗi và phương án dự phòng bằng luật đảm nhiệm vai trò ổn định hóa quy trình. Vì vậy, phương pháp~1 không nên được hiểu như một thí nghiệm lời nhắc đơn lẻ, mà là một pipeline AI/data có khả năng chạy trên dữ liệu lớn, kiểm tra lại được và mở rộng được.

Trong bài toán tóm tắt đánh giá khách sạn, một yêu cầu quan trọng là tránh làm mất các ý kiến nhỏ nhưng có giá trị. Nếu một khách sạn có rất nhiều đánh giá tích cực về phòng nhưng cũng có một nhóm phàn nàn lặp lại về tiếng ồn hoặc bữa sáng, tóm tắt cần phản ánh được cả điểm mạnh lẫn rủi ro. Nếu chỉ dùng cảm xúc tổng thể, phản hồi tiêu cực thiểu số dễ bị che khuất. Nếu chỉ dùng tóm tắt tự do, mô hình có thể viết trôi chảy nhưng không chỉ ra bằng chứng cụ thể. Phương pháp~1 giải quyết vấn đề này bằng cách giữ lại các đoạn bằng chứng trung gian, sau đó gom chúng theo khía cạnh, cảm xúc và cụm chủ đề trước khi sinh tóm tắt.

\subsection{Mục tiêu thiết kế}

Mục tiêu thiết kế đầu tiên là bảo đảm tính bám nguồn. Mỗi đơn vị đầu ra của quy trình cần có khả năng truy ngược về đánh giá gốc hoặc ít nhất là đoạn văn bản nguồn hỗ trợ cho nhận định đó. Trong triển khai, các trường như \texttt{data\_source}, \texttt{hotel\_id}, \texttt{review\_index}, \texttt{source\_review}, \texttt{source\_text} và \texttt{segment\_text} được dùng để giữ mối liên hệ này. Nhờ vậy, khi tóm tắt cuối nêu rằng khách hàng khen vị trí thuận tiện hoặc phàn nàn về tiếng ồn, người phân tích có thể lần ngược về các đoạn đánh giá cụ thể đã tạo ra kết luận đó. Đây là cơ sở để gọi phương pháp này là tóm tắt có căn cứ bằng chứng, thay vì chỉ là tóm tắt sinh tự do.

Mục tiêu thiết kế thứ hai là phân rã bài toán thành các bước đủ nhỏ để kiểm tra được. Một đánh giá khách sạn thường không phù hợp với cách xử lý một nhãn duy nhất, vì nó có thể chứa nhiều mục tiêu và nhiều chiều cảm xúc. Do đó, quy trình tách nhiệm vụ lớn thành các bước: tiền phân đoạn, trích xuất đoạn ABSA, kiểm soát khía cạnh, phân loại cảm xúc, gán cụm chủ đề, gom bằng chứng và sinh tóm tắt. Mỗi bước có đầu vào, đầu ra và điều kiện hợp lệ riêng. Khi một tóm tắt cuối có vấn đề, hệ thống không cần suy đoán chung chung rằng mô hình sai, mà có thể kiểm tra xem lỗi xuất hiện ở bước tách đơn vị ý kiến, gán khía cạnh, phân loại cảm xúc, gán cụm hay viết tóm tắt.

Mục tiêu thiết kế thứ ba là giữ ổn định khi gọi API trên quy mô lớn. Trong bối cảnh thực nghiệm nhỏ, việc gửi một số đoạn văn cho LLM có thể khá đơn giản. Tuy nhiên, khi số lượng đánh giá tăng lên, các vấn đề kỹ thuật trở nên quan trọng: giới hạn tốc độ, hết thời gian chờ, lỗi mạng, phản hồi không đúng định dạng, chi phí gọi API và khả năng phải chạy lại sau khi dừng giữa chừng. Vì vậy, phương pháp được thiết kế với bộ nhớ đệm, điểm lưu tiến độ, cơ chế thử lại có giới hạn và phương án dự phòng bằng luật. Những cơ chế này không trực tiếp cải thiện ngữ nghĩa của mô hình, nhưng quyết định việc quy trình có thể vận hành ổn định hay không.

Mục tiêu thiết kế thứ tư là tạo ra dữ liệu trung gian phục vụ đánh giá. Một tóm tắt cuối dù đọc tốt vẫn chưa đủ để chứng minh quy trình đúng. Cần kiểm tra được các lớp trước đó: đơn vị ý kiến có bám nguồn không, nhãn khía cạnh có đúng bộ phân loại không, cảm xúc có đúng với mục tiêu đang xét không, cụm chủ đề có phản ánh nhóm bằng chứng thật không và tóm tắt có bao phủ các cụm quan trọng không. Vì vậy, quy trình tạo các tệp trung gian như dữ liệu câu đã xử lý, đầu ra theo khía cạnh, cụm bằng chứng và các bảng chỉ số. Các tệp này giúp việc đánh giá không dừng ở cảm nhận về văn bản cuối.

\subsection{Luồng xử lý tổng quan}

Ở mức tổng quát, phương pháp~1 có thể được mô tả như một chuỗi xử lý từ dữ liệu thô đến tóm tắt có cấu trúc. Đầu vào là các tệp CSV chứa mã đánh giá và nội dung đánh giá. Hệ thống đọc dữ liệu theo lô, phân tích mã đánh giá để lấy nguồn dữ liệu, mã khách sạn và chỉ số đánh giá, sau đó làm sạch nội dung văn bản. Các đánh giá hợp lệ được chuyển thành các đơn vị xử lý nhỏ hơn. Trên từng đơn vị này, LLM thực hiện trích xuất ABSA để tạo các đoạn có khía cạnh và văn bản chuẩn hóa. Sau đó hệ thống hậu kiểm khía cạnh, phân loại cảm xúc, gán cụm, gom bằng chứng và sinh tóm tắt.

\begin{figure}[H]
\centering
\fbox{
\begin{minipage}{0.94\textwidth}
\centering
\vspace{0.35cm}
\begin{tabular}{c}
\textbf{Tệp CSV đánh giá khách sạn} \\
$\downarrow$ \\
\textbf{Đọc theo lô và kiểm tra định danh} \\
\small{Phân tích \texttt{ref\_id}, lấy nguồn dữ liệu, khách sạn và chỉ số đánh giá} \\
$\downarrow$ \\
\textbf{Làm sạch văn bản và tạo khóa truy vết} \\
\small{Bỏ đánh giá rỗng, chuẩn hóa văn bản, giữ lại đánh giá gốc} \\
$\downarrow$ \\
\textbf{Tiền phân đoạn} \\
\small{Tách đánh giá thành câu hoặc đơn vị ý kiến bám nguồn} \\
$\downarrow$ \\
\textbf{Trích xuất ABSA bằng LLM/API} \\
\small{Tạo đoạn ý kiến có khía cạnh, bằng chứng, ngôn ngữ và độ tin cậy} \\
$\downarrow$ \\
\textbf{Kiểm soát khía cạnh} \\
\small{Chuẩn hóa hoặc chuyển nhãn theo bộ phân loại sáu khía cạnh} \\
$\downarrow$ \\
\textbf{Phân loại cảm xúc} \\
\small{Gán tích cực, tiêu cực hoặc trung tính cho từng đoạn theo khía cạnh} \\
$\downarrow$ \\
\textbf{Gán cụm chủ đề} \\
\small{Gán mã cụm, nhãn cụm và mô tả cho từng đoạn bằng chứng} \\
$\downarrow$ \\
\textbf{Gom bằng chứng} \\
\small{Gom theo khách sạn, khía cạnh, cảm xúc và cụm chủ đề} \\
$\downarrow$ \\
\textbf{Sinh tóm tắt theo cụm bằng chứng}
\end{tabular}
\vspace{0.35cm}
\end{minipage}
}
\caption{Luồng xử lý chính của quy trình LLM/API trong phương pháp 1}
\label{fig:method1-workflow}
\end{figure}

Sơ đồ trong Hình~\ref{fig:method1-workflow} cố ý chỉ thể hiện các bước nghiệp vụ chính. Trong triển khai thực tế, mỗi bước có gọi API còn được bao quanh bởi các lớp vận hành như bộ nhớ đệm, thử lại, kiểm tra phản hồi và ghi trạng thái. Các lớp này không làm thay đổi bản chất nghiệp vụ của phương pháp, nhưng giúp quy trình đủ ổn định để chạy trên nhiều tệp đánh giá. Khi viết báo cáo, cần phân biệt rõ hai loại thành phần: các bước phân tích ngôn ngữ tạo ra kết quả nghiên cứu và các cơ chế vận hành giúp quá trình đó có thể chạy lại, kiểm tra lại và mở rộng.

\subsection{Đọc dữ liệu, lọc đầu vào và định danh truy vết}

Giai đoạn đầu tiên của phương pháp là đọc dữ liệu và thiết lập định danh. Đầu vào của quy trình là các tệp CSV đánh giá khách sạn, trong đó mỗi dòng chứa mã tham chiếu và nội dung đánh giá. Mã tham chiếu được phân tích để lấy thông tin về nguồn dữ liệu, mã khách sạn và chỉ số đánh giá. Từ các thành phần này, hệ thống tạo khóa khách sạn ổn định để gom các đánh giá cùng khách sạn trong các bước sau. Việc tạo khóa ngay từ đầu giúp tránh tình trạng tóm tắt bị lẫn giữa các khách sạn hoặc không thể truy ngược kết quả về nguồn.

Dữ liệu được đọc theo lô thay vì nạp toàn bộ vào bộ nhớ. Đây là lựa chọn quan trọng khi xử lý dữ liệu đánh giá ở quy mô lớn. Đọc theo lô giúp giảm tiêu thụ bộ nhớ, cho phép ghi tiến độ sau từng phần và hỗ trợ chạy lại khi tác vụ bị dừng. Ngoài ra, quy trình có thể áp dụng các điều kiện lọc như giới hạn số dòng, giới hạn số khách sạn hoặc chỉ xử lý các khách sạn có số đánh giá tối thiểu. Các tùy chọn này không làm thay đổi thuật toán chính, nhưng giúp việc chạy thử, kiểm tra mẫu và đánh giá trên tập lớn linh hoạt hơn.

Trong giai đoạn này, các bản ghi không hợp lệ được loại bỏ sớm. Nếu mã tham chiếu không đúng định dạng, nếu văn bản đánh giá rỗng hoặc nếu khách sạn không thỏa điều kiện lọc, bản ghi đó không được đưa vào các bước gọi LLM. Việc lọc sớm có hai lợi ích. Thứ nhất, nó tránh lãng phí chi phí API cho các dữ liệu không sử dụng được. Thứ hai, nó giúp thống kê vận hành rõ ràng hơn, vì hệ thống có thể ghi lại số dòng đã đọc, số dòng lỗi định danh, số đánh giá rỗng và số dòng bị bỏ qua do điều kiện lọc.

Định danh truy vết cũng được dùng trong các đầu ra trung gian. Khi một đoạn ABSA được tạo ra, nó vẫn giữ thông tin về khách sạn, đánh giá gốc, câu hoặc đơn vị ý kiến ban đầu và đoạn bằng chứng. Nhờ vậy, các bước tổng hợp sau có thể gom chính xác theo khách sạn, đồng thời người phân tích có thể kiểm tra lại từng kết luận. Đây là một yêu cầu thực tế quan trọng: nếu tóm tắt không thể truy vết về bản ghi nguồn, việc đánh giá độ tin cậy của tóm tắt sẽ phụ thuộc quá nhiều vào cảm nhận chủ quan.

\subsection{Tiền phân đoạn đánh giá}

Tiền phân đoạn là bước chuẩn bị văn bản trước khi trích xuất ABSA. Mục tiêu của bước này là chia đánh giá thành các đơn vị đủ nhỏ để mỗi đơn vị có thể được phân tích khía cạnh và cảm xúc một cách rõ ràng. Đơn vị xử lý có thể là một câu được tách theo dấu câu hoặc một đơn vị ý kiến được LLM chuẩn hóa từ đánh giá gốc. Trong cả hai trường hợp, hệ thống đều cố gắng giữ mối liên hệ với văn bản nguồn để tránh biến tiền phân đoạn thành tóm tắt tự do.

Trong dữ liệu thực tế, đánh giá khách sạn thường không được viết theo cấu trúc chuẩn. Người dùng có thể viết câu rất dài, thiếu dấu câu, trộn nhiều ngôn ngữ, lặp từ, dùng biểu tượng cảm xúc hoặc nối nhiều nhận xét trong cùng một dòng. Nếu đưa trực tiếp toàn bộ đánh giá này vào bước ABSA, mô hình dễ bỏ sót một số mục tiêu hoặc gán cảm xúc theo ấn tượng chung. Tiền phân đoạn giúp giảm vấn đề đó bằng cách tách hoặc chuẩn hóa đánh giá thành các đơn vị tự đủ ngữ cảnh, trong đó mỗi đơn vị vẫn giữ mục tiêu và ý kiến chính.

Khi dùng LLM cho tiền phân đoạn, lời nhắc yêu cầu mô hình chỉ làm nhiệm vụ chuẩn hóa đơn vị ý kiến, không gán khía cạnh, không gán cảm xúc và không sinh kết luận mới. Đầu ra mong muốn gồm văn bản đơn vị, đoạn nguồn hỗ trợ, độ tin cậy và một số chỉ báo phụ nếu có. Ranh giới này rất quan trọng: nếu mô hình bắt đầu diễn giải, tóm tắt hoặc thêm ý mới ngay ở bước tiền phân đoạn, các bước sau sẽ xử lý trên dữ liệu đã bị biến đổi quá xa nguồn. Vì vậy, bước tiền phân đoạn phải được hiểu là làm rõ và chia nhỏ văn bản, không phải là bước ra quyết định nội dung.

Chất lượng tiền phân đoạn ảnh hưởng trực tiếp đến toàn bộ phần sau của quy trình. Nếu một ý kiến bị bỏ mất ở đây, các bước khía cạnh, cảm xúc và tóm tắt không thể khôi phục lại nó. Nếu một đơn vị bị tách quá nhỏ, ví dụ chỉ còn một tính từ mà không có mục tiêu, bước ABSA có thể không biết tính từ đó nói về phòng, nhân viên hay bữa sáng. Nếu một đơn vị bị gộp quá lớn, nhiều khía cạnh và cảm xúc sẽ bị trộn lẫn. Do đó, tiền phân đoạn là một bước đáng chú trọng trong phương pháp, không chỉ là tiền xử lý kỹ thuật.

\subsection{Trích xuất đoạn ABSA}

Sau tiền phân đoạn, từng câu hoặc đơn vị ý kiến được đưa vào bước trích xuất ABSA. Bước này yêu cầu LLM tách văn bản thành các đoạn ý kiến tập trung, trong đó mỗi đoạn có một mục tiêu chính và một khía cạnh chính. Một đơn vị đầu vào có thể tạo ra không đoạn nào nếu nó là nhiễu, một đoạn nếu nó chứa một ý kiến rõ, hoặc nhiều đoạn nếu nó chứa nhiều khía cạnh khác nhau. Đây là điểm quan trọng: quy trình không ép mỗi câu hoặc mỗi đánh giá về một nhãn duy nhất.

Đầu ra của bước trích xuất gồm các trường như \texttt{segment\_text}, \texttt{aspect}, \texttt{detected\_language}, \texttt{normalized\_text\_vi}, \texttt{normalized\_text\_en}, \texttt{split\_reason} và \texttt{confidence}. Trường \texttt{segment\_text} biểu diễn đoạn ý kiến đã được tách; \texttt{aspect} là khía cạnh do LLM đề xuất; các trường chuẩn hóa tiếng Việt và tiếng Anh giúp các bước sau chọn ngôn ngữ phù hợp cho phân loại và tóm tắt. Độ tin cậy không được xem là thước đo tuyệt đối, nhưng là tín hiệu hữu ích khi kiểm tra hoặc phân tích lỗi.

Bộ khía cạnh trong phương pháp gồm sáu nhóm. \texttt{facility} bao gồm cơ sở vật chất và môi trường vật lý như phòng, giường, phòng tắm, view, vị trí, tiếng ồn, nội thất và điều hòa. \texttt{amenity} bao gồm tiện ích khách sạn cung cấp như WiFi, bữa sáng, hồ bơi, nhà hàng, spa, phòng gym, bãi đỗ xe hoặc dịch vụ đưa đón. \texttt{service} bao gồm tương tác con người và quy trình phục vụ như lễ tân, nhân viên, nhận phòng, trả phòng, hỗ trợ và xử lý vấn đề. \texttt{experience} phản ánh cảm nhận tổng thể, sự thoải mái, thư giãn và giá trị trải nghiệm. \texttt{branding} liên quan đến uy tín, chuẩn thương hiệu hoặc hình ảnh cao cấp. \texttt{loyalty} liên quan đến ý định quay lại, giới thiệu hoặc gắn bó.

Việc trích xuất ở mức đoạn giúp xử lý các câu nhiều khía cạnh tốt hơn. Ví dụ, câu ``Phòng sạch nhưng bữa sáng ít lựa chọn và nhân viên rất thân thiện'' không nên có một nhãn duy nhất. Quy trình cần tạo ít nhất ba đoạn: một đoạn tích cực về cơ sở vật chất, một đoạn tiêu cực về tiện ích ăn uống và một đoạn tích cực về dịch vụ. Nhờ vậy, khi tổng hợp ở cuối, điểm mạnh và điểm yếu của khách sạn không bị trộn lẫn.

\subsection{Kiểm soát khía cạnh bằng luật}

Sau khi LLM đề xuất nhãn khía cạnh cho từng đoạn, hệ thống áp dụng một lớp kiểm soát bằng luật nhằm chuẩn hóa nhãn về bộ phân loại sáu nhóm cố định và loại bỏ các nhãn không hợp lệ. Lớp này cần thiết vì LLM có thể đề xuất các nhãn ngoài taxonomy, dùng từ đồng nghĩa hoặc viết không nhất quán giữa các lần gọi. Nếu không kiểm soát, bước gom bằng chứng theo khía cạnh sẽ bị phân tán: cùng một nhóm ý kiến về ``nhân viên'' có thể rơi vào ba nhãn khác nhau như \texttt{service}, \texttt{staff} và \texttt{front\_desk}, làm mất khả năng tổng hợp.

Cơ chế kiểm soát gồm ba thành phần. Thứ nhất, một bộ ánh xạ từ đồng nghĩa (synonym mapping) chuyển các nhãn biến thể về nhãn chuẩn: \texttt{staff}, \texttt{front\_desk}, \texttt{reception}, \texttt{clerk} đều được đưa về \texttt{service}; \texttt{room}, \texttt{bed}, \texttt{bathroom}, \texttt{ac}, \texttt{aircon} được đưa về \texttt{facility}; \texttt{breakfast}, \texttt{wifi}, \texttt{pool}, \texttt{gym}, \texttt{parking} được đưa về \texttt{amenity}; \texttt{location}, \texttt{view}, \texttt{noise}, \texttt{neighborhood} được đưa về \texttt{facility}; \texttt{value}, \texttt{overall}, \texttt{experience}, \texttt{comfort} được đưa về \texttt{experience}; \texttt{brand}, \texttt{reputation}, \texttt{luxury} được đưa về \texttt{branding}; \texttt{return}, \texttt{recommend}, \texttt{loyalty} được đưa về \texttt{loyalty}.

Thứ hai, một bộ luật ưu tiên xử lý trường hợp một đoạn được gán nhiều khía cạnh. Nếu LLM trả về danh sách nhãn, luật chọn nhãn có độ tin cậy cao nhất làm khía cạnh chính và ghi các nhãn còn lại làm khía cạnh phụ. Nếu không có thông tin độ tin cậy, luật ưu tiên theo thứ tự: \texttt{facility} $>$ \texttt{amenity} $>$ \texttt{service} $>$ \texttt{experience} $>$ \texttt{branding} $>$ \texttt{loyalty}, dựa trên tần suất xuất hiện trong dữ liệu đánh giá khách sạn.

Thứ ba, các nhãn không thuộc bộ phân loại sáu nhóm được đánh dấu là \texttt{unknown} và không đưa vào bước gom bằng chứng. Việc loại bỏ sớm các nhãn ngoài taxonomy giúp tránh tình trạng một nhóm bằng chứng ``rác'' xuất hiện trong đầu ra cuối. Tuy nhiên, các đoạn mang nhãn \texttt{unknown} vẫn được lưu lại trong tệp trung gian để phục vụ phân tích lỗi và kiểm tra độ phủ của bộ phân loại.

Lớp kiểm soát bằng luật không thay thế năng lực hiểu ngữ nghĩa của LLM, mà đóng vai trò ổn định hóa nhãn ở quy mô lớn. Khi chạy trên hàng nghìn đánh giá, việc dựa hoàn toàn vào LLM để gán nhãn sẽ tạo ra sự không nhất quán do biến thể từ vựng và ngẫu nhiên của mô hình. Bộ luật ánh xạ và ưu tuyến giúp đảm bảo rằng cùng một loại ý kiến được gom về cùng một nhóm, từ đó bước tổng hợp bằng chứng và sinh tóm tắt hoạt động trên dữ liệu nhất quán.

\section{Phương pháp 2: SemAE huấn luyện trên SPACE và suy luận trên HASOS}

Phương pháp~2 sử dụng SemAE làm bộ chọn bằng chứng theo khía cạnh. Mô hình được huấn luyện trên SPACE, sau đó dùng checkpoint đã huấn luyện để suy luận trên HASOS~\cite{hasos}. HASOS có 50 khách sạn, 5.000 đánh giá, 45.529 câu và taxonomy 29 khía cạnh. Các khía cạnh con được gom về bốn nhóm có tóm tắt tham chiếu (gold summary) khi tính ROUGE: Facility, Amenity, Service và Experience.

\subsection{Luồng chọn bằng chứng}

Đầu tiên, hệ thống tách đánh giá thành câu. Mỗi câu được mã hóa bằng bộ mã hóa (encoder) của SemAE. Với mỗi khía cạnh trong taxonomy HASOS, hệ thống xây dựng prototype từ bộ từ khóa hạt giống (seed words) tương ứng. Câu đánh giá được xếp hạng theo độ phù hợp với prototype, sau đó lọc bằng ngưỡng điểm và giới hạn token đầu vào. Các câu được chọn trở thành bằng chứng cho bước sinh hoặc ghép tóm tắt.

\begin{figure}[tbp]
\centering
\begin{adjustbox}{max width=\textwidth}
\begin{tikzpicture}[node distance=0.55cm and 0.55cm]
    \node[smallblock, text width=2.25cm] (reviews) {Đánh giá HASOS};
    \node[smallblock, text width=2.0cm, right=of reviews] (sent) {Tách câu};
    \node[smallblock, text width=2.2cm, right=of sent] (enc) {Mã hóa SemAE};
    \node[smallblock, text width=2.45cm, right=of enc] (proto) {Prototype\\29 khía cạnh};
    \node[smallblock, text width=2.35cm, right=of proto] (rank) {Xếp hạng câu\\theo khía cạnh};
    \node[smallblock, text width=2.15cm, right=of rank] (pool) {Pool bằng chứng\\$score \leq T$};

    \node[smallblock, text width=2.2cm, below=0.8cm of pool] (m2) {M2\\rewrite trực tiếp};
    \node[smallblock, text width=2.35cm, below=0.55cm of m2] (m3) {M3\\keyword split};
    \node[smallblock, text width=2.35cm, below=0.55cm of m3] (m4) {M4\\BERT-ABSA split};
    \node[smallblock, text width=2.2cm, above=0.55cm of m2] (m1) {M1\\extractive};

    \node[smallblock, text width=2.7cm, right=0.75cm of m1] (out1) {Câu gốc\\cắt $B$ từ};
    \node[smallblock, text width=2.7cm, right=0.75cm of m2] (out2) {FLAN-T5\\không tách cảm xúc};
    \node[smallblock, text width=2.7cm, right=0.75cm of m3] (out3) {FLAN-T5\\pos/neg theo từ khóa};
    \node[smallblock, text width=2.7cm, right=0.75cm of m4] (out4) {FLAN-T5\\pos/neg theo BERT};

    \draw[arrow] (reviews) -- (sent);
    \draw[arrow] (sent) -- (enc);
    \draw[arrow] (enc) -- (proto);
    \draw[arrow] (proto) -- (rank);
    \draw[arrow] (rank) -- (pool);
    \draw[arrow] (pool) |- (m1);
    \draw[arrow] (pool) -- (m2);
    \draw[arrow] (pool) |- (m3);
    \draw[arrow] (pool) |- (m4);
    \draw[arrow] (m1) -- (out1);
    \draw[arrow] (m2) -- (out2);
    \draw[arrow] (m3) -- (out3);
    \draw[arrow] (m4) -- (out4);
\end{tikzpicture}
\end{adjustbox}
\caption{Workflow phương pháp 2 và điểm rẽ nhánh của các biến thể M1--M4}
\label{fig:semae-hasos-pipeline}
\end{figure}

Hình~\ref{fig:semae-hasos-pipeline} tách rõ hai phần của phương pháp~2. Phần bên trái là pipeline chung: tách câu, mã hóa bằng SemAE, tạo prototype cho 29 khía cạnh và xếp hạng câu để lấy pool bằng chứng. Phần bên phải là nơi bốn biến thể khác nhau: M1 giữ nguyên câu gốc, M2 sinh tóm tắt từ pool chưa tách cảm xúc, M3 tách cảm xúc bằng từ khóa trước khi sinh, còn M4 thay bước tách cảm xúc bằng backend BERT-ABSA.

\subsection{Các biến thể M1--M4}

Để tách riêng tác động của từng thành phần, hệ thống triển khai bốn biến thể. Điểm chung của cả bốn biến thể là đều xuất phát từ bằng chứng do SemAE xếp hạng theo khía cạnh. Điểm khác nhau nằm ở ba quyết định hậu xử lý: có sinh tóm tắt hay không, có tách cảm xúc trước khi sinh hay không, và nếu có thì dùng backend nào để tách cảm xúc.

\begin{itemize}[leftmargin=1.2cm]
    \item \textbf{M1}: giữ nguyên các câu SemAE chọn, đóng vai trò baseline trích rút (extractive). M1 không dùng decoder sinh văn bản, không có nhánh sinh thành công/thất bại và không tách cảm xúc.
    \item \textbf{M2}: dùng bằng chứng đã chọn để sinh tóm tắt trừu tượng (abstractive) thông qua mô hình FLAN-T5, chưa tách cảm xúc. M2 được triển khai ở chế độ phân cấp (hierarchical): tóm tắt cấp khía cạnh con $\to$ tóm tắt cấp khía cạnh cha.
    \item \textbf{M3}: tách bằng chứng theo cảm xúc bằng backend từ khóa (keyword) rồi sinh tóm tắt riêng cho từng nhóm cảm xúc dương (pos) và âm (neg).
    \item \textbf{M4}: tách bằng chứng theo cảm xúc bằng backend BERT-ABSA rồi sinh tóm tắt riêng cho từng nhóm cảm xúc. BERT-ABSA dùng mô hình dựa trên BERT để gán nhãn cảm xúc ở cấp câu, mang lại nhãn chính xác hơn backend từ khóa.
\end{itemize}

\begin{table}[tbp]
\centering
\caption{So sánh workflow của các biến thể M1--M4}
\label{tab:m1-m4-workflow}
\small
\begin{tabularx}{\textwidth}{l L{0.24\textwidth} L{0.23\textwidth} Y}
\toprule
\textbf{Biến thể} & \textbf{Bằng chứng đầu vào} & \textbf{Cách tạo đầu ra} & \textbf{Mục đích so sánh} \\
\midrule
M1 & Câu SemAE đã chọn, không tách cảm xúc. & Ghép/cắt câu gốc theo ngân sách $B$; không sinh văn bản mới. & Baseline trích rút để đo mức trung thực khi không có hallucination từ decoder. \\
M2 & Pool bằng chứng theo khía cạnh, chưa tách cảm xúc. & Sinh tóm tắt trừu tượng trực tiếp từ pool bằng chứng bằng FLAN-T5. & Đo tác động của bước sinh tóm tắt khi chưa kiểm soát riêng cảm xúc. \\
M3 & Bằng chứng đã tách pos/neg bằng luật và từ khóa cảm xúc. & Sinh tóm tắt riêng cho các nhóm cảm xúc hợp lệ. & Kiểm tra liệu tách cảm xúc đơn giản có cải thiện ROUGE và giảm lẫn cực tính hay không. \\
M4 & Bằng chứng đã tách pos/neg bằng backend BERT-ABSA. & Sinh tóm tắt riêng cho các nhóm cảm xúc do mô hình ABSA xác định. & Kiểm tra lợi ích của tách cảm xúc theo mô hình học máy so với luật/từ khóa. \\
\bottomrule
\end{tabularx}
\end{table}

Vì vậy, M1 không nên được diễn giải như một hệ sinh tóm tắt cùng loại với M2--M4. M1 là điểm neo trích rút để kiểm tra faithfulness; M2--M4 mới là các biến thể có bước sinh tóm tắt. Trong nhóm sinh, M2 là baseline không tách cảm xúc, M3 kiểm tra tách cảm xúc bằng luật/từ khóa, còn M4 kiểm tra tách cảm xúc bằng mô hình BERT-ABSA.

\subsection{Cơ chế trung thực bằng chứng}

Một rủi ro cơ bản của sinh tóm tắt trừu tượng là hallucination: mô hình sinh có thể đưa vào tóm tắt những thông tin không có trong bằng chứng, đặc biệt khi pool bằng chứng lớn và mô hình có xu hướng diễn giải tự do. Để giảm rủi ro này, các biến thể sinh M2--M4 được thiết kế với một cơ chế trung thực bằng chứng gồm năm thành phần. Cơ chế này không thay thế mô hình sinh mà đóng vai trò lớp kiểm soát hậu xử lý, đảm bảo mọi câu trong tóm tắt cuối đều có cơ sở trong bằng chứng được chọn.

\textbf{Thành phần 1 -- Giới hạn bằng chứng đầu vào.} Số câu bằng chứng được đưa vào mô hình sinh bị giới hạn ở top-$k$ câu có điểm SemAE cao nhất, với $k=5$. Giới hạn này khớp với số bằng chứng mà LLM judge sử dụng khi chấm ($judge\_evidence\_k = 5$), đảm bảo rằng mọi thông tin mô hình thấy khi sinh đều có mặt trong tập bằng chứng judge kiểm tra. Nếu không giới hạn, mô hình sinh có thể dựa vào câu bằng chứng xếp hạng thấp (thứ 6--18) để sinh ra thông tin mà judge không thể xác nhận, dẫn đến cờ lỗi ``khẳng định không căn cứ''.

\textbf{Thành phần 2 -- Lời nhắc ràng buộc trung thực.} Lời nhắc (prompt) gửi cho mô hình sinh được bổ sung các ràng buộc rõ ràng: mọi事实 trong tóm tắt phải là diễn đạt lại (paraphrase) của ít nhất một câu bằng chứng; nếu chi tiết không có trong bằng chứng thì bỏ qua; nếu bằng chứng mỏng thì viết tóm tắt ngắn thay vì điền thêm đoán; giữ nguyên cực tính (polarity) của bằng chứng, không lật khen thành chê hoặc ngược lại. Lời nhắc còn kèm một ví dụ minh họa về cách viết tóm tắt trung thực từ bằng chứng, giúp mô hình hiểu ranh giới giữa diễn đạt lại và bịa thêm.

\textbf{Thành phần 3 -- Bộ lọc câu không được hỗ trợ.} Sau khi mô hình sinh tóm tắt, mỗi câu trong tóm tắt được kiểm tra độ trùng lặp nội dung (content-word overlap) với từng câu bằng chứng. Độ trùng lặp được tính trên tập từ nội dung (loại bỏ stopword và chuẩn hóa số nhiều) bằng hệ số overlap coefficient: tỷ lệ giữa phần giao và tập từ nhỏ hơn. Nếu câu sinh có độ trùng lặp cao nhất với mọi câu bằng chứng đều dưới ngưỡng 0{,}30, câu đó được xem là không được hỗ trợ và bị thay thế bằng câu bằng chứng có thứ hạng cao nhất chưa được dùng. Bộ lọc này loại bỏ hallucination ở cấp câu mà không cần đọc lại toàn bộ tóm tắt.

\textbf{Thành phần 4 -- Kiểm tra nhất quán cảm xúc.} Đối với M3 và M4, hai biến thể tách cảm xúc, một rủi ro đặc biệt là lật cảm xúc (sentiment flip): mô hình sinh có thể viết tóm tắt tích cực cho bucket tiêu cực hoặc ngược lại do nhiễu trong nhãn cảm xúc đầu vào. Để giảm rủi ro, cực tính của tóm tắt sinh ra được ước lượng bằng một bộ từ điển cảm xúc (lexicon) đơn giản: nếu tóm tắt có nhiều từ tích cực hơn tiêu cực, cực tính là ``pos''; ngược lại là ``neg''. Nếu cực tính của tóm tắt mâu thuẫn với bucket mục tiêu, hệ thống không xuất bản tóm tắt đó mà chuyển sang cơ chế dự phòng. Ngoài ra, khi thay thế câu không được hỗ trợ (thành phần 3), pool bằng chứng thay thế cũng được lọc theo cực tính: chỉ câu bằng chứng có cực tính trùng với bucket mục tiêu (hoặc trung tính) mới được dùng, đảm bảo cơ chế thay thế không đưa thêm câu trái cực tính vào tóm tắt.

\textbf{Thành phần 5 -- Dự phòng khử trùng lặp.} Khi mô hình sinh thất bại (đầu ra rỗng, lặp lại prompt, hoặc quá chung chung), hệ thống chuyển sang cơ chế dự phòng: ghép các câu bằng chứng top-rank thành tóm tắt. Tuy nhiên, pool bằng chứng có thể chứa các câu gần trùng lặp (cùng ý, khác cách diễn đạt), làm tóm tắt dự phòng bị lặp ý. Để khắc phục, các câu bằng chứng được khử trùng lặp trước khi ghép: mỗi câu mới được so sánh với các câu đã giữ bằng độ trùng lặp nội dung, và nếu độ trùng lặp vượt ngưỡng 0{,}82, câu đó bị bỏ. Cơ chế này giúp tóm tắt dự phòng ngắn gọn và không lặp ý, đồng thời giữ nguyên tính bám nguồn.

Năm thành phần trên hoạt động phối hợp: thành phần~1 và~2 giảm hallucination ở bước sinh, thành phần~3 loại bỏ hallucination còn sót ở bước hậu xử lý, thành phần~4 kiểm soát cực tính cho biến thể tách cảm xúc, và thành phần~5 đảm bảo cơ chế dự phòng không tạo ra tóm tắt kém chất lượng. Cơ chế này áp dụng chung cho M2--M4, nhưng tác động mạnh nhất lên M3 và M4 vì hai biến thể này có thêm rủi ro lật cảm xúc từ bước tách bucket.

\subsection{Vai trò của các tham số $T$ và $B$}

Hai tham số $T$ và $B$ không có cùng ý nghĩa tuyệt đối ở cả bốn biến thể. Tham số $T$ là ngưỡng chọn bằng chứng của SemAE. Vì điểm SemAE trong pipeline được diễn giải như độ lệch/khoảng cách, câu có điểm nhỏ hơn hoặc bằng $T$ được giữ lại trong pool bằng chứng. Tham số $B$ kiểm soát độ dài đầu ra: với M1, $B$ là giới hạn số từ của chuỗi câu trích rút; với M2--M4, $B$ là \texttt{max\_new\_tokens} của mô hình sinh tóm tắt.

M1 là baseline trích rút từ run \texttt{space\_hasos\_full\_e20}. Biến thể này không dùng decoder sinh văn bản, không có nhánh sinh thành công/thất bại, do đó các cột \textit{Gen.} và \textit{Fallback} không áp dụng. M1 dùng $B=40$ từ theo cấu hình baseline chính thức. Để kiểm tra câu hỏi liệu $B$ của M1 có nên lớn hơn hay không, báo cáo chạy thêm một sweep hậu nghiệm trên log bằng chứng dài hơn \texttt{space\_hasos\_threshold\_full}: $B \in \{40,64,80,96,120\}$. Kết quả trong Bảng~\ref{tab:m1-token-sweep} cho thấy $B=40$ vẫn tốt nhất theo ROUGE-1 trong replay này, vì tăng độ dài làm đưa thêm câu ít khớp tham chiếu vào tóm tắt trích rút.

M2--M4 được sweep theo cả $T$ và $B$ vì cùng pool bằng chứng có thể được lọc lại rồi sinh lại tóm tắt mà không cần huấn luyện lại SemAE. Thiết lập cuối cùng được dùng trong so sánh là M2 với $T=0{,}0075, B=128$, M3 với $T=0{,}0055, B=96$, và M4 với $T=0{,}005, B=96$.

\section{Thiết kế đánh giá}

Đánh giá được chia thành ba lớp. Lớp thứ nhất là ROUGE-1/2/L, dùng để so khớp với tóm tắt tham chiếu (gold summary). Lớp thứ hai là các chỉ số tự động từ artifact như tỷ lệ dự phòng (fallback rate), tỷ lệ sao chép bằng chứng (copy rate), số bằng chứng trung bình và tỷ lệ nén (compression ratio). Lớp thứ ba là bộ chấm tự động bằng mô hình ngôn ngữ lớn (LLM judge) để chấm độ trung thực, đúng khía cạnh, đúng cảm xúc và mức độ hữu ích của bản tóm tắt.

\begin{figure}[tbp]
\centering
\begin{tikzpicture}[node distance=0.8cm and 0.9cm]
    \node[block, text width=2.6cm] (outputs) {Đầu ra M1--M4};
    \node[block, right=of outputs, text width=2.7cm] (rouge) {ROUGE\\tóm tắt tham chiếu};
    \node[block, above=of rouge, text width=2.7cm] (judge) {LLM judge\\rubric JSON};
    \node[block, below=of rouge, text width=2.7cm] (auto) {Chỉ số artifact\\dự phòng, sao chép, nén};
    \node[block, right=of rouge, text width=3.0cm] (tables) {Bảng kết quả\\so sánh M1--M4};
    \draw[arrow] (outputs) -- (rouge);
    \draw[arrow] (outputs) |- (judge);
    \draw[arrow] (outputs) |- (auto);
    \draw[arrow] (rouge) -- (tables);
    \draw[arrow] (judge) -| (tables);
    \draw[arrow] (auto) -| (tables);
\end{tikzpicture}
\caption{Luồng đánh giá kết hợp ROUGE, chỉ số tự động và LLM judge}
\label{fig:evaluation-flow}
\end{figure}

Hình~\ref{fig:evaluation-flow} cho thấy mỗi nhóm chỉ số trả lời một câu hỏi khác nhau: ROUGE đo mức trùng khớp với tóm tắt tham chiếu, chỉ số artifact đo đặc tính vận hành, còn LLM judge đo độ trung thực và đúng khía cạnh.

Một item được tính là đạt chuẩn trong LLM judge nếu thỏa các điều kiện: điểm hỗ trợ bằng chứng, đúng khía cạnh và đúng cảm xúc đều từ 4 trở lên; độ bao phủ (coverage) từ 3 trở lên; không có khẳng định không căn cứ (unsupported claim) và không có lỗi lật cảm xúc (sentiment flip). Do đó, tỷ lệ vượt chuẩn (pass rate) không phải là ROUGE, mà là tỷ lệ tóm tắt vượt qua ngưỡng chất lượng ngữ nghĩa tối thiểu theo bộ tiêu chí chấm.

% =========================
% CHUONG 4
% =========================
\chapter{Thực nghiệm và đánh giá}

\section{Thiết lập thực nghiệm}

Thực nghiệm phương pháp~2 sử dụng checkpoint SemAE huấn luyện trên SPACE và suy luận trên HASOS. SPACE được dùng như tập đối chiếu vì có các khía cạnh tổng quát và tóm tắt tham chiếu (gold summary) theo định dạng gần với thiết kế SemAE ban đầu. HASOS là tập mục tiêu chính, có taxonomy 29 khía cạnh và dữ liệu đánh giá khách sạn ở cấp thực thể.

Trong HASOS, vì tóm tắt tham chiếu (gold summary) chỉ bao phủ bốn nhóm khía cạnh cha, các kết quả ở 29 khía cạnh con được gom về bốn nhóm khi tính ROUGE. Hai nhóm Branding và Loyalty không được đưa vào mẫu số ROUGE của HASOS vì không có tóm tắt tham chiếu tương ứng. Cách xử lý này tránh việc báo điểm trên những trường hợp không có tham chiếu hợp lệ.

\begin{table}[tbp]
\centering
\caption{Tóm tắt dữ liệu HASOS sử dụng trong phương pháp 2}
\label{tab:hasos-data-summary}
\begin{tabularx}{\textwidth}{L{0.28\textwidth} r Y}
\toprule
\textbf{Thành phần} & \textbf{Số lượng} & \textbf{Ghi chú} \\
\midrule
Khách sạn / thực thể & 50 & Mỗi thực thể gồm nhiều đánh giá người dùng. \\
Đánh giá & 5.000 & Dữ liệu đầu vào dạng văn bản tự nhiên. \\
Câu đánh giá & 45.529 & Đơn vị được SemAE xếp hạng làm bằng chứng. \\
Khía cạnh HASOS & 29 & Taxonomy chi tiết theo miền khách sạn. \\
Nhóm khía cạnh có tóm tắt tham chiếu & 4 & Facility, Amenity, Service, Experience. \\
\bottomrule
\end{tabularx}
\end{table}

\begin{table}[tbp]
\centering
\caption{Thiết lập sử dụng trong so sánh cuối cùng trên HASOS}
\label{tab:optimal-settings}
\begin{threeparttable}
\begin{tabularx}{\textwidth}{l c c Y}
\toprule
\textbf{Biến thể} & \textbf{Ngưỡng $T$} & \textbf{Ngân sách $B$} & \textbf{Diễn giải} \\
\midrule
M1 & -- & 40 từ & Baseline trích rút SemAE từ \texttt{space\_hasos\_full\_e20}; không dùng decoder sinh và không có fallback. \\
M2 & 0.0075 & 128 token & Tóm tắt trừu tượng trước khi tách cảm xúc. \\
M3 & 0.0055 & 96 token & Tách cảm xúc bằng luật/từ khóa, sau đó sinh tóm tắt. \\
M4 & 0.0050 & 96 token & Tách cảm xúc bằng BERT-ABSA, sau đó sinh tóm tắt. \\
\bottomrule
\end{tabularx}
\begin{tablenotes}
\footnotesize
\item Với M1, $B$ là giới hạn số từ của tóm tắt trích rút. Với M2--M4, $B$ là \texttt{max\_new\_tokens}; $T$ là ngưỡng giữ bằng chứng $score \leq T$.
\end{tablenotes}
\end{threeparttable}
\end{table}

\begin{table}[tbp]
\centering
\caption{Độ nhạy ngân sách trích rút $B$ của M1}
\label{tab:m1-token-sweep}
\small
\begin{tabular}{r r r r}
\toprule
\textbf{$B$ (từ)} & \textbf{ROUGE-1} & \textbf{ROUGE-2} & \textbf{ROUGE-L} \\
\midrule
40  & \textbf{\num{0.2078}} & \textbf{\num{0.0591}} & \textbf{\num{0.1245}} \\
64  & \num{0.1592} & \num{0.0524} & \num{0.0986} \\
80  & \num{0.1374} & \num{0.0479} & \num{0.0869} \\
96  & \num{0.1210} & \num{0.0448} & \num{0.0781} \\
120 & \num{0.1031} & \num{0.0403} & \num{0.0680} \\
\bottomrule
\end{tabular}

\vspace{0.35em}
\begin{minipage}{0.82\textwidth}
\footnotesize\RaggedRight
\textit{Ghi chú.} Đây là kiểm tra độ nhạy hậu nghiệm, không phải một phương pháp mới. Replay dùng log bằng chứng dài hơn của run \texttt{space\_hasos\_threshold\_full} để kiểm tra $B \leq 120$; kết quả ROUGE chính của M1 trong các bảng so sánh vẫn lấy từ baseline \texttt{space\_hasos\_full\_e20} với $B=40$.
\end{minipage}
\end{table}

Các bảng trong phần này được tổ chức theo từng nhóm câu hỏi thực nghiệm. Bảng~\ref{tab:hasos-data-summary} mô tả dữ liệu, Bảng~\ref{tab:optimal-settings} và Bảng~\ref{tab:m1-token-sweep} ghi cấu hình so sánh và kiểm tra độ nhạy của M1. Bảng~\ref{tab:evidence-faithfulness} và Bảng~\ref{tab:judge-errors} tách riêng chất lượng bằng chứng và lỗi judge. Bảng~\ref{tab:summary-utility} và Bảng~\ref{tab:summary-quality} trình bày chất lượng tóm tắt. Bảng~\ref{tab:production-ops}, Bảng~\ref{tab:production-rouge}, Bảng~\ref{tab:space-rouge}, Bảng~\ref{tab:metric-availability} và Bảng~\ref{tab:reporting-principles} gom các chỉ số vận hành, ROUGE, phạm vi hợp lệ và nguyên tắc diễn giải kết quả.

Cần lưu ý rằng \textit{N} trong các bảng judge không phải là cùng một mẫu số cố định cho mọi biến thể, mà là số item tóm tắt hợp lệ thực sự được đưa vào chấm. M1 có 1.370 item vì baseline \texttt{space\_hasos\_full\_e20} chỉ sinh output không rỗng cho 1.370 cặp thực thể--khía cạnh con. M2 có 1.372 item từ pool threshold tối ưu. M3 và M4 có ít item hơn (907 và 804) vì hai biến thể này tách bằng chứng theo cảm xúc dương/âm; các ô không có bằng chứng hoặc không có output hợp lệ không được đưa vào judge. M4 còn dùng backend BERT-ABSA nghiêm ngặt hơn nên số ô hợp lệ thấp hơn M3. Vì vậy, \textit{N} phản ánh phạm vi output hợp lệ của từng pipeline, không phải điểm chất lượng.

\FloatBarrier

\section{Kết quả sweep tham số}

Sweep tham số cho thấy SPACE giữ ngưỡng mặc định $T=0{,}0082$, trong khi HASOS cần thiết lập riêng theo phương pháp. Như Bảng~\ref{tab:optimal-settings} trình bày, M2 cải thiện khi tăng ngưỡng lên 0{,}0075, M3 cải thiện mạnh ở 0{,}0055, còn M4 giữ ngưỡng mặc định 0{,}005. Token budget tốt nhất cho ba biến thể trừu tượng lần lượt là 128, 96 và 96. Các thiết lập này được dùng làm base cho các bảng chỉ số cuối cùng.

\section{Kết quả ROUGE}

\begin{table}[tbp]
\centering
\caption{ROUGE trên HASOS}
\label{tab:production-rouge}
\begin{tabular}{l r r r}
\toprule
\textbf{Phương pháp} & \textbf{ROUGE-1} & \textbf{ROUGE-2} & \textbf{ROUGE-L} \\
\midrule
M1 & \num{0.203} & \num{0.055} & \num{0.122} \\
M2 & \num{0.232} & \num{0.044} & \num{0.150} \\
M3 & \textbf{\num{0.262}} & \textbf{\num{0.071}} & \textbf{\num{0.168}} \\
M4 & \num{0.209} & \num{0.040} & \num{0.137} \\
\bottomrule
\end{tabular}
\end{table}

Theo ROUGE trên HASOS ở Bảng~\ref{tab:production-rouge}, M3 đạt ROUGE-1 cao nhất với \num{0.262}, tiếp theo là M2 với \num{0.232} và M4 với \num{0.209}. Điều này cho thấy tách cảm xúc bằng từ khóa (keyword) trước khi sinh tóm tắt giúp tăng độ trùng khớp với tóm tắt tham chiếu. Tuy nhiên, ROUGE chỉ phản ánh độ trùng n-gram; nó không khẳng định bản tóm tắt có hoàn toàn trung thực với bằng chứng hay không.

\begin{table}[tbp]
\centering
\caption{ROUGE trên SPACE dùng như phụ lục đối chiếu}
\label{tab:space-rouge}
\begin{tabularx}{\textwidth}{Y r r r}
\toprule
\textbf{Phương pháp} & \textbf{ROUGE-1} & \textbf{ROUGE-2} & \textbf{ROUGE-L} \\
\midrule
M1 trích rút (extractive) & \num{0.304} & \num{0.087} & \num{0.222} \\
M2 trừu tượng (abstractive) & \num{0.307} & \num{0.086} & \num{0.235} \\
M3 tách cảm xúc bằng từ khóa & \textbf{\num{0.309}} & \textbf{\num{0.089}} & \textbf{\num{0.240}} \\
M4 tách cảm xúc bằng BERT-ABSA & \num{0.308} & \textbf{\num{0.089}} & \num{0.239} \\
\bottomrule
\end{tabularx}
\end{table}

Kết quả SPACE trong Bảng~\ref{tab:space-rouge} cho thấy bốn biến thể có điểm ROUGE khá gần nhau, với M3 nhỉnh hơn ở ROUGE-1 và ROUGE-L. Khoảng cách trên HASOS rõ hơn vì HASOS dùng taxonomy chi tiết hơn, làm cho việc chọn và tổ chức bằng chứng theo khía cạnh có ảnh hưởng lớn hơn.

\section{Đánh giá bằng LLM judge}

LLM judge được chạy trên 4.453 item HASOS bằng một mô hình ngôn ngữ lớn, với bộ tiêu chí chấm (rubric) cố định, đầu ra JSON được kiểm tra hợp lệ và có cache để tái lập. Mỗi item được chấm theo các trục hỗ trợ bằng chứng, đúng khía cạnh, đúng cảm xúc, độ bao phủ (coverage), độ cụ thể (specificity) và độ hữu ích (usefulness), đồng thời có các cờ lỗi như khẳng định không căn cứ (unsupported claim), sai khía cạnh và lật cảm xúc (sentiment flip). Các chỉ số này được tách thành Bảng~\ref{tab:evidence-faithfulness}, Bảng~\ref{tab:judge-errors}, Bảng~\ref{tab:summary-utility} và Bảng~\ref{tab:summary-quality} để dễ đọc hơn.

\begin{table}[tbp]
\centering
\caption{Chất lượng bằng chứng và độ trung thực theo LLM judge}
\label{tab:evidence-faithfulness}
\begin{threeparttable}
\begin{tabular}{l r r r r}
\toprule
\textbf{Phương pháp} & \textbf{N} & \textbf{P@5} & \textbf{Support@5} & \makecell{\textbf{Điểm hỗ trợ}\\\textbf{trung bình}} \\
\midrule
M1 & 1370 & \num{0.860} & \num{0.876} & \num{4.401} \\
M2 & 1372 & \num{0.855} & \num{0.695} & \num{4.353} \\
M3 &  907 & \num{0.909} & \num{0.841} & \num{4.530} \\
M4 &  804 & \textbf{\num{0.926}} & \textbf{\num{0.875}} & \textbf{\num{4.627}} \\
\bottomrule
\end{tabular}
\begin{tablenotes}
\footnotesize
\item \textit{N} = số item hợp lệ được đưa vào judge; \textit{P@5} = tỷ lệ bằng chứng đúng trong năm câu đầu; \textit{Support@5} = tỷ lệ bằng chứng đầu hỗ trợ nhận định; điểm hỗ trợ trung bình theo thang 1--5.
\end{tablenotes}
\end{threeparttable}
\end{table}

\begin{table}[tbp]
\centering
\caption{Các tỷ lệ lỗi theo LLM judge}
\label{tab:judge-errors}
\begin{threeparttable}
\begin{tabular}{l r r r r}
\toprule
\textbf{Phương pháp} & \makecell{\textbf{Leak}\\\textbf{khía cạnh}} & \makecell{\textbf{Leak}\\\textbf{cảm xúc}} & \makecell{\textbf{Không}\\\textbf{căn cứ}} & \makecell{\textbf{Lật}\\\textbf{cảm xúc}} \\
\midrule
M1 & \num{0.067} & \textbf{\num{0.006}} & \textbf{\num{0.093}} & \textbf{\num{0.016}} \\
M2 & \textbf{\num{0.044}} & \textbf{\num{0.006}} & \num{0.147} & \num{0.020} \\
M3 & \num{0.053} & \num{0.024} & \num{0.110} & \num{0.099} \\
M4 & \num{0.050} & \num{0.013} & \num{0.091} & \num{0.061} \\
\bottomrule
\end{tabular}
\begin{tablenotes}
\footnotesize
\item Các cột đều là tỷ lệ lỗi nên giá trị thấp hơn tốt hơn. M4 giảm tốt nhất cả leak khía cạnh lẫn khẳng định không căn cứ trong nhóm biến thể sinh.
\end{tablenotes}
\end{threeparttable}
\end{table}

Bảng~\ref{tab:evidence-faithfulness} cho thấy M4 đạt P@5 cao nhất (\num{0.926}), Support@5 cao nhất (\num{0.875}) và điểm hỗ trợ trung bình cao nhất (\num{4.627}) trong cả bốn biến thể. M3 đứng thứ hai về P@5 (\num{0.909}) và điểm hỗ trợ (\num{4.530}). M1, dù là baseline trích rút, chỉ đạt P@5 = \num{0.860} và Support@5 = \num{0.876}; điều này cho thấy các biến thể sinh với cơ chế trung thực bằng chứng không chỉ ngang bằng mà còn vượt M1 về chất lượng bằng chứng được judge nhận diện.

Bảng~\ref{tab:judge-errors} cho thấy M4 có tỷ lệ khẳng định không căn cứ thấp nhất (\num{0.091}), thấp hơn cả M1 (\num{0.093}). M2 giảm khẳng định không căn cứ xuống \num{0.147} nhờ cơ chế trung thực, trong khi không có cơ chế thì tỷ lệ này cao hơn đáng kể. Lật cảm xúc ở M4 (\num{0.061}) thấp hơn M3 (\num{0.099}) nhờ backend BERT-ABSA gán nhãn cảm xúc chính xác hơn backend từ khóa.

\begin{table}[tbp]
\centering
\caption{Đánh giá đúng khía cạnh và cảm xúc theo LLM judge}
\label{tab:summary-utility}
\begin{threeparttable}
\begin{tabular}{l r r r r}
\toprule
\textbf{Phương pháp} & \makecell{\textbf{Tỷ lệ}\\\textbf{vượt chuẩn}} & \makecell{\textbf{Đúng}\\\textbf{khía cạnh}} & \makecell{\textbf{Đúng}\\\textbf{cảm xúc}} & \makecell{\textbf{Nhớ}\\\textbf{chủ đề}} \\
\midrule
M1 & \num{0.720} & \num{4.456} & \num{4.593} & \num{0.903} \\
M2 & \num{0.700} & \num{4.491} & \num{4.522} & \num{0.880} \\
M3 & \num{0.734} & \num{4.527} & \num{4.512} & \num{0.916} \\
M4 & \textbf{\num{0.787}} & \textbf{\num{4.568}} & \textbf{\num{4.673}} & \textbf{\num{0.917}} \\
\bottomrule
\end{tabular}
\begin{tablenotes}
\footnotesize
\item Các điểm đúng khía cạnh và đúng cảm xúc theo thang 1--5; tỷ lệ vượt chuẩn và nhớ chủ đề nằm trong khoảng 0--1.
\end{tablenotes}
\end{threeparttable}
\end{table}

Bảng~\ref{tab:summary-utility} là kết quả trọng tâm của báo cáo. Trên HASOS, M4 đạt tỷ lệ vượt chuẩn cao nhất với \num{0.787}, tiếp theo là M3 với \num{0.734}, M1 với \num{0.720} và M2 với \num{0.700}. Thứ tự này cho thấy các biến thể sinh có tách cảm xúc (M3, M4) đã vượt baseline trích rút M1 theo tiêu chí ngữ nghĩa của LLM judge. M4 cũng đạt điểm đúng khía cạnh cao nhất (\num{4.568}), điểm đúng cảm xúc cao nhất (\num{4.673}) và tỷ lệ nhớ chủ đề cao nhất (\num{0.917}). Kết quả này khẳng định rằng cơ chế trung thực bằng chứng kết hợp với backend BERT-ABSA giúp M4 tạo ra tóm tắt vừa trung thực vừa đúng cảm xúc, đồng thời bao phủ chủ đề tốt hơn baseline trích rút.

\begin{table}[tbp]
\centering
\caption{Đánh giá độ bao phủ và tính hữu ích theo LLM judge}
\label{tab:summary-quality}
\begin{threeparttable}
\begin{tabular}{l r r r r}
\toprule
\textbf{Phương pháp} & \makecell{\textbf{Bao}\\\textbf{phủ}} & \makecell{\textbf{Cụ}\\\textbf{thể}} & \makecell{\textbf{Hữu}\\\textbf{ích}} & \makecell{\textbf{Chung}\\\textbf{chung}} \\
\midrule
M1 & \num{4.002} & \num{4.067} & \num{4.043} & \textbf{\num{0.096}} \\
M2 & \num{3.804} & \num{3.910} & \num{3.985} & \num{0.149} \\
M3 & \num{4.166} & \num{4.205} & \num{4.179} & \num{0.110} \\
M4 & \textbf{\num{4.300}} & \textbf{\num{4.333}} & \textbf{\num{4.310}} & \num{0.100} \\
\bottomrule
\end{tabular}
\begin{tablenotes}
\footnotesize
\item Bao phủ, cụ thể và hữu ích theo thang 1--5. Cột chung chung là tỷ lệ lỗi nên giá trị thấp hơn tốt hơn.
\end{tablenotes}
\end{threeparttable}
\end{table}

Bảng~\ref{tab:summary-quality} cho thấy M4 dẫn đầu về độ bao phủ (\num{4.300}), độ cụ thể (\num{4.333}) và độ hữu ích (\num{4.310}). M3 đứng thứ hai trong nhóm sinh. M1 có tỷ lệ tóm tắt chung chung thấp nhất (\num{0.096}), nhưng các điểm bao phủ và cụ thể thấp hơn M3--M4, phù hợp với bản chất trích rút: M1 giữ nguyên câu bằng chứng nên không bị chung chung, nhưng cũng không có khả năng tổng hợp và diễn đạt lại.

\begin{table}[tbp]
\centering
\caption{Chỉ số vận hành trên HASOS}
\label{tab:production-ops}
\begin{threeparttable}
\begin{tabular}{l r c c r r r}
\toprule
\textbf{Phương pháp} & \textbf{Rows} & \textbf{Gen.} & \textbf{Fallback} & \textbf{Copy} & \textbf{Ev.} & \textbf{Comp.} \\
\midrule
M1 & 1370 & -- & -- & \num{1.000} & \num{2.613} & \num{1.000} \\
M2 & 1372 & \num{0.562} & \num{0.329} & \num{0.185} & \num{18.389} & \num{0.491} \\
M3 &  907 & \num{0.428} & \num{0.436} & \num{0.521} & \num{2.834} & \num{0.732} \\
M4 &  804 & \num{0.345} & \num{0.515} & \num{0.614} & \num{1.919} & \num{0.822} \\
\bottomrule
\end{tabular}
\begin{tablenotes}
\footnotesize
\item \textit{Gen.} và \textit{Fallback} không áp dụng cho M1 vì M1 không sinh văn bản; M1 chỉ cắt và ghép câu bằng chứng đã chọn. \textit{Ev.} = số bằng chứng trung bình; \textit{Comp.} = tỷ lệ nén. \textit{Fallback} của M3/M4 bao gồm cả dự phòng do lật cảm xúc và do đầu ra chung chung.
\end{tablenotes}
\end{threeparttable}
\end{table}

Bảng~\ref{tab:production-ops} cho thấy M2 có tỷ lệ sinh thành công cao nhất (\num{0.562}) và tỷ lệ sao chép thấp nhất (\num{0.185}) trong nhóm sinh, nhưng số bằng chứng trung bình rất cao (\num{18.389}) do không tách cảm xúc. M3 và M4 có tỷ lệ dự phòng cao hơn vì cơ chế trung thực bằng chứng chuyển sang dự phòng khi phát hiện lật cảm xúc hoặc đầu ra chung chung. Tỷ lệ sao chép của M4 (\num{0.614}) cao hơn M2, phản ánh việc cơ chế dự phòng ghép câu bằng chứng khi mô hình sinh thất bại.

\begin{table}[tbp]
\centering
\caption{Phạm vi hợp lệ của các nhóm chỉ số}
\label{tab:metric-availability}
\begin{tabularx}{\textwidth}{L{0.23\textwidth} L{0.20\textwidth} Y}
\toprule
\textbf{Nhóm chỉ số} & \textbf{Trạng thái} & \textbf{Diễn giải} \\
\midrule
ROUGE-1/2/L & Tính trực tiếp & Dùng tóm tắt tham chiếu theo nhóm khía cạnh; HASOS được gom 29 khía cạnh con về 4 nhóm có tham chiếu. \\
Bằng chứng và độ trung thực & Dựa trên judge & LLM judge chấm 4.453 item M1--M4 theo rubric cố định, có cache và schema JSON. \\
Tỷ lệ nén / dự phòng / sao chép & Tính trực tiếp & Lấy từ artifact sinh tóm tắt và tệp bằng chứng JSONL; Gen./Fallback không áp dụng cho M1. \\
ABSA Macro-F1 / Pair-F1 & Chưa báo chính thức & Dữ liệu hiện tại không có nhãn gold khía cạnh--cảm xúc ở cấp câu, nên không giả lập F1 như chỉ số gold. \\
Quad exact/partial F1 & Chưa báo chính thức & Cần nhãn gold ở cấp span, hiện không tồn tại trong artifact. \\
p95 latency / cache / retry & Chưa báo chính thức & Cần instrumentation runtime chi tiết; báo cáo hiện chỉ có chi phí judge và các tỷ lệ từ artifact. \\
\bottomrule
\end{tabularx}
\end{table}

\begin{table}[tbp]
\centering
\caption{Nguyên tắc diễn giải kết quả thực nghiệm}
\label{tab:reporting-principles}
\small
\begin{tabularx}{\textwidth}{L{0.24\textwidth} Y Y}
\toprule
\textbf{Vấn đề} & \textbf{Nên trình bày trong báo cáo} & \textbf{Không nên diễn giải} \\
\midrule
ROUGE & Chỉ số so khớp n-gram giữa tóm tắt hệ thống và tóm tắt tham chiếu; phù hợp để so sánh độ gần bề mặt văn bản. & Không xem ROUGE là thước đo đầy đủ của độ đúng ngữ nghĩa hoặc độ trung thực. \\
LLM judge & Chỉ số dựa trên rubric để ước lượng độ hỗ trợ bằng chứng, đúng khía cạnh, đúng cảm xúc và tính hữu ích. & Không gọi pass rate là accuracy theo nghĩa gold label hoặc thay thế đánh giá của chuyên gia người. \\
M1 & Baseline trích rút: giữ câu bằng chứng do SemAE chọn và cắt theo ngân sách từ. & Không ghi M1 có bước sinh, fallback, hallucination từ decoder hoặc sweep tối ưu giống M2--M4. \\
M2--M4 & Nhóm biến thể sinh tóm tắt có cơ chế trung thực bằng chứng: M2 chưa tách cảm xúc, M3 tách bằng luật/từ khóa, M4 tách bằng BERT-ABSA. & Không so M1 với M2--M4 như các decoder cùng bản chất; M1 chủ yếu là mốc đối chứng về trích rút. \\
\textit{N} khác nhau & \textit{N} là số item hợp lệ sau khi từng pipeline tạo output; M3/M4 ít hơn do tách cảm xúc làm mất các ô không có bằng chứng hợp lệ. & Không xem \textit{N} lớn/nhỏ là điểm chất lượng trực tiếp, và không nói mọi phương pháp được judge trên cùng mẫu số. \\
Metric thiếu gold & Nói rõ ABSA F1 và quad F1 chưa báo chính thức vì thiếu nhãn gold ở cấp câu/span. & Không tự tạo hoặc suy diễn F1 từ judge, ROUGE hoặc output của hệ thống. \\
\bottomrule
\end{tabularx}
\end{table}

\section{Phân tích tổng hợp}

Hai nhóm chỉ số đưa ra hai góc nhìn khác nhau. Nếu ưu tiên giống tóm tắt tham chiếu theo ROUGE, M3 là lựa chọn tốt nhất trên HASOS với ROUGE-1 = \num{0.262}. Nếu ưu tiên độ trung thực và đúng cảm xúc theo LLM judge, M4 là lựa chọn mạnh nhất với pass rate = \num{0.787}, vượt cả baseline trích rút M1 (\num{0.720}). Thứ tự pass rate trên HASOS là M4 $>$ M3 $>$ M1 $>$ M2, cho thấy các biến thể sinh có tách cảm xúc đã vượt baseline trích rút nhờ cơ chế trung thực bằng chứng kết hợp backend cảm xúc chất lượng cao.

Kết quả này cho thấy cơ chế trung thực bằng chứng đã thay đổi câu chuyện so sánh giữa trích rút và sinh trừu tượng. Trước khi áp dụng cơ chế này, baseline trích rút M1 thường đạt pass rate cao nhất vì không có hallucination từ decoder. Sau khi áp dụng, các biến thể sinh M3 và M4 giảm khẳng định không căn cứ và lật cảm xúc đủ để vượt M1, đồng thời giữ được lợi thế về độ bao phủ, độ cụ thể và khả năng diễn đạt lại.

Trên SPACE (Bảng~\ref{tab:space-rouge}), bốn biến thể có ROUGE gần nhau, cho thấy khi taxonomy đơn giản hơn (6 khía cạnh), sự khác biệt giữa các biến thể nhỏ hơn. Kết quả LLM judge trên SPACE cho thấy M3 và M4 cải thiện và vượt M2, nhưng vẫn nằm gần M1. Điều này phù hợp với kỳ vọng: SPACE có ít khía cạnh hơn, pool bằng chứng gọn hơn, nên lợi thế của tách cảm xúc và cơ chế trung thực ít rõ hơn trên HASOS.

Từ kết quả trên, có thể rút ra một hướng dẫn thực tiễn. Với dashboard phân tích nội bộ cần độ tin cậy và tránh lật cảm xúc, M4 là lựa chọn phù hợp nhất. Với báo cáo theo tóm tắt tham chiếu hoặc benchmark truyền thống, M3 có thể được ưu tiên do điểm ROUGE tốt hơn. Một hướng kết hợp hợp lý là dùng M3 để tối ưu độ trùng khớp và dùng LLM judge như lớp kiểm soát chất lượng trước khi xuất bản tóm tắt.

\section{Hạn chế}

Báo cáo không trình bày ABSA Macro-F1 hoặc Quad F1 như chỉ số chính thức vì dữ liệu hiện tại không có nhãn gold khía cạnh--cảm xúc ở cấp câu hoặc span. Nếu tự dùng nhãn sinh bởi hệ thống để chấm lại hệ thống, kết quả sẽ không có giá trị như đánh giá trên nhãn gold. Ngoài ra, các chỉ số vận hành như p95 latency, tỷ lệ thử lại (retry rate) hoặc tỷ lệ trúng cache (cache hit rate) của pipeline summarization gốc cần instrumentation runtime chi tiết hơn trước khi báo cáo.

Một hạn chế khác là LLM judge vẫn phụ thuộc vào bộ tiêu chí chấm (rubric) và mô hình đánh giá. Dù toàn bộ output được kiểm tra hợp lệ bằng schema và có cache để tái lập, chỉ số dựa trên judge không thay thế hoàn toàn đánh giá thủ công. Trong phiên bản tiếp theo, có thể chọn một tập mẫu nhỏ để người đánh giá kiểm tra chéo nhằm ước lượng độ nhất quán giữa LLM judge và đánh giá của con người.

% =========================
% KET LUAN
% =========================
\chapter{Kết luận}

Khóa luận đã trình bày bài toán phân tích cảm xúc theo khía cạnh và tóm tắt đánh giá khách sạn theo hai hướng tiếp cận. Phương pháp thứ nhất dựa trên LLM/API, phù hợp với mục tiêu tạo đầu ra linh hoạt và có cấu trúc. Phương pháp thứ hai sử dụng SemAE huấn luyện trên SPACE và suy luận trên HASOS, giúp xây dựng một baseline có tính tái lập cao hơn ở bước chọn bằng chứng.

Kết quả thực nghiệm trên HASOS cho thấy các biến thể sinh M3 và M4, được trang bị cơ chế trung thực bằng chứng và backend tách cảm xúc, đã vượt baseline trích rút M1 theo tỷ lệ vượt chuẩn của LLM judge: M4 đạt \num{0.787}, M3 đạt \num{0.734}, M1 đạt \num{0.720}. Về ROUGE, M3 đạt ROUGE-1 cao nhất (\num{0.262}), cho thấy tách cảm xúc bằng từ khóa giúp tăng độ trùng khớp với tóm tắt tham chiếu. Kết quả này khẳng định cần đánh giá hệ thống tóm tắt bằng nhiều nhóm chỉ số: ROUGE đo độ trùng khớp bề mặt, còn LLM judge đo độ trung thực, đúng cảm xúc và tính hữu ích.

Về mặt kỹ thuật, cơ chế trung thực bằng chứng gồm năm thành phần (giới hạn bằng chứng đầu vào, lời nhắc ràng buộc, bộ lọc câu không được hỗ trợ, kiểm tra nhất quán cảm xúc và dự phòng khử trùng lặp) đã giảm hallucination đủ để các biến thể sinh vượt baseline trích rút. Backend BERT-ABSA của M4 gán nhãn cảm xúc chính xác hơn backend từ khóa của M3, giúp M4 giảm lật cảm xúc và đạt pass rate cao nhất.

Trong tương lai, đề tài có thể được mở rộng theo ba hướng. Thứ nhất, xây dựng bộ nhãn gold ở cấp câu hoặc span để đánh giá ABSA Macro-F1, Pair-F1 và Quad F1 một cách chính thức. Thứ hai, bổ sung instrumentation vận hành để đo p95 latency, tỷ lệ thử lại (retry rate), tỷ lệ trúng cache (cache hit rate) và chi phí trên mỗi 1.000 đánh giá. Thứ ba, thay backend tách cảm xúc của M3 bằng mô hình phân tích cảm xúc dựa trên Transformer (ví dụ \texttt{cardiffnlp/twitter-roberta-base-sentiment}) để nâng cao chất lượng bucket cảm xúc, đồng thời mở rộng cơ chế trung thực bằng chứng sang các miền dữ liệu ngoài khách sạn.

\begin{thebibliography}{99}

\bibitem{hu2004mining}
Minqing Hu and Bing Liu.
\newblock Mining and summarizing customer reviews.
\newblock In \textit{Proceedings of the 10th ACM SIGKDD International Conference on Knowledge Discovery and Data Mining}, pages 168--177, 2004.
\newblock \url{https://doi.org/10.1145/1014052.1014073}

\bibitem{pontiki2014semeval}
Maria Pontiki, Dimitris Galanis, John Pavlopoulos, Harris Papageorgiou, Ion Androutsopoulos, and Suresh Manandhar.
\newblock SemEval-2014 Task 4: Aspect Based Sentiment Analysis.
\newblock In \textit{Proceedings of the 8th International Workshop on Semantic Evaluation (SemEval 2014)}, pages 27--35, 2014.
\newblock \url{https://aclanthology.org/S14-2004/}

\bibitem{lin2004rouge}
Chin-Yew Lin.
\newblock ROUGE: A Package for Automatic Evaluation of Summaries.
\newblock In \textit{Text Summarization Branches Out}, 2004.
\newblock \url{https://aclanthology.org/W04-1013/}

\bibitem{vaswani2017attention}
Ashish Vaswani et al.
\newblock Attention Is All You Need.
\newblock In \textit{NeurIPS}, 2017.
\newblock \url{https://arxiv.org/abs/1706.03762}

\bibitem{devlin2019bert}
Jacob Devlin, Ming-Wei Chang, Kenton Lee, and Kristina Toutanova.
\newblock BERT: Pre-training of Deep Bidirectional Transformers for Language Understanding.
\newblock In \textit{Proceedings of NAACL-HLT 2019}, pages 4171--4186, 2019.
\newblock \url{https://aclanthology.org/N19-1423/}

\bibitem{brown2020language}
Tom B. Brown et al.
\newblock Language Models are Few-Shot Learners.
\newblock In \textit{Advances in Neural Information Processing Systems}, 2020.
\newblock \url{https://arxiv.org/abs/2005.14165}

\bibitem{angelidis2021extractive}
Stefanos Angelidis, Reinald Kim Amplayo, Yoshihiko Suhara, Xiaolan Wang, and Mirella Lapata.
\newblock Extractive Opinion Summarization in Quantized Transformer Spaces.
\newblock \textit{Transactions of the Association for Computational Linguistics}, 9:277--293, 2021.
\newblock \url{https://aclanthology.org/2021.tacl-1.17/}

\bibitem{basuroychowdhury2022sparse}
Somnath Basu Roy Chowdhury, Chao Zhao, and Snigdha Chaturvedi.
\newblock Unsupervised Extractive Opinion Summarization Using Sparse Coding.
\newblock In \textit{Proceedings of the 60th Annual Meeting of the Association for Computational Linguistics}, pages 1209--1225, 2022.
\newblock \url{https://aclanthology.org/2022.acl-long.86/}

\bibitem{bhaskar2023prompted}
Adithya Bhaskar, Alex Fabbri, and Greg Durrett.
\newblock Prompted Opinion Summarization with GPT-3.5.
\newblock In \textit{Findings of the Association for Computational Linguistics: ACL 2023}, pages 9282--9300, 2023.
\newblock \url{https://aclanthology.org/2023.findings-acl.591/}

\bibitem{liu2023geval}
Yang Liu, Dan Iter, Yichong Xu, Shuohang Wang, Ruochen Xu, and Chenguang Zhu.
\newblock G-Eval: NLG Evaluation using GPT-4 with Better Human Alignment.
\newblock In \textit{Proceedings of the 2023 Conference on Empirical Methods in Natural Language Processing}, pages 2511--2522, 2023.
\newblock \url{https://aclanthology.org/2023.emnlp-main.153/}

\bibitem{zheng2023judging}
Lianmin Zheng et al.
\newblock Judging LLM-as-a-Judge with MT-Bench and Chatbot Arena.
\newblock \textit{arXiv preprint arXiv:2306.05685}, 2023.
\newblock \url{https://arxiv.org/abs/2306.05685}

\bibitem{hasos}
HASOS hotel review benchmark artifacts.
\newblock Local dataset and taxonomy: \texttt{data/hasos/hasos\_summ.json}, \texttt{data/hasos/aspect\_taxonomy.json}, and \texttt{data/seeds\_hasos/}.

\end{thebibliography}

\end{document}
